% ==============================================================================
% Fichier     : Cours_Securite.tex
% Titre       : Sécurité Informatique
% Auteur      : MINKA MI NGUIDJOÏ Thierry Emmanuel
%               Sr-Eng, Sr-Auditor, GRC Expert
% Institution : Département de Génie Informatique
%               École Nationale Supérieure Polytechnique (ENSPY)
%               Université de Yaoundé~I
% Niveau      : Humanités Numériques, Niveau~3 (Licence 3)
% Année       : 2025--2026
% Description : Fichier principal du cours. Orchestre tous les chapitres,
%               la bibliographie et les éléments de couverture.
% Compilateur : pdflatex + biber (ou lualatex recommandé)
%               → lualatex Cours && biber Cours && lualatex Cours && lualatex Cours
% ==============================================================================

\documentclass[12pt, a4paper, twoside, openright]{book}

% ==============================================================================
% CHARGEMENT DE LA FEUILLE DE STYLE
% ==============================================================================
\usepackage{Style_cours}

% ==============================================================================
% BIBLIOGRAPHIE
% ==============================================================================
\addbibresource{bibliographie.bib}

% ==============================================================================
% MÉTADONNÉES PDF
% ==============================================================================
\hypersetup{
    pdftitle  = {Sécurité Informatique},
    pdfauthor = {MINKA MI NGUIDJOÏ Thierry Emmanuel},
    pdfsubject= {Support de Cours— Sécurité SI, Cybersécurité, Gestion des Risques},
    pdfkeywords={Sécurité SI, Cybersécurité, Risk Management, Cryptographie,
                 Sécurité réseau, Université de Yaoundé I}
}

% ==============================================================================
% DÉBUT DU DOCUMENT
% ==============================================================================
\begin{document}

% ------------------------------------------------------------------------------
% 1. PAGE DE TITRE (1ère DE COUVERTURE)
% ------------------------------------------------------------------------------
\begin{titlepage}
\thispagestyle{empty}
\pagecolor{CouleurPrimaire}

\begin{tikzpicture}[remember picture, overlay]
    % Bande dorée en bas de couverture
    \fill[CouleurAccent]
        (current page.south west) rectangle
        ([yshift=2.2cm] current page.south east);
    % Ligne décorative en haut
    \fill[CouleurAccent]
        ([yshift=-1.5cm] current page.north west) rectangle
        ([yshift=-1.8cm] current page.north east);
    % Élément graphique décoratif (carré tourné)
    \fill[white, opacity=0.04]
        ([xshift=8cm, yshift=-6cm] current page.north west)
        -- ++(5,0) -- ++(0,-5) -- ++(-5,0) -- cycle;
    \fill[white, opacity=0.06]
        ([xshift=9cm, yshift=-7cm] current page.north west)
        -- ++(6,0) -- ++(0,-6) -- ++(-6,0) -- cycle;
\end{tikzpicture}

\vspace*{1.5cm}

% Logo institutionnel (à décommenter si disponible)
% \begin{center}
%     \includegraphics[height=2.5cm]{figures/logo_uy1.png}
% \end{center}

\begin{center}
    {\color{white!80!gray}\small\scshape
        Université de Yaoundé~I~·~Ecole Nationale Supérieure Polytechnique~·~Département de Génie Informatique}\\[0.4ex]
    {\color{CouleurAccent}\rule{0.55\linewidth}{0.8pt}}
\end{center}

\vspace{1.5cm}

\begin{center}
    {\color{white!60!gray}\footnotesize\scshape
        support de cours}\\[1ex]
    {\color{white}\fontsize{34}{42}\selectfont\bfseries
        Sécurité Informatique}\\[1.5ex]
    {\color{CouleurAccent}\Large\itshape
        Fondamentaux, Pratiques et Gouvernance de la Cybersécurité}
\end{center}

\vspace{2cm}

\begin{center}
    {\color{white!70!gray}\small
        \ding{43}\quad Humanités Numériques — Niveau 3 (Licence 3)}\\[0.5ex]
    {\color{white!70!gray}\small
        \ding{44}\quad Security and Risk Management ~·~ Cryptographie ~·~ Sécurité des Réseaux}\\[0.5ex]
    {\color{white!70!gray}\small
        \ding{228}\quad Gestion des Identités ~·~ Tests d'Intrusion ~·~ Sécurité des Applications}
\end{center}

\vfill

\begin{center}
    {\color{CouleurAccent}\rule{0.7\linewidth}{1.2pt}}\\[1.2ex]
    {\color{white}\large\bfseries MINKA MI NGUIDJOÏ Thierry Emmanuel}\\[0.4ex]
    {\color{white!80!gray}\small
        CISA~·~CISM~·~CGEIT~·~CRISC~·~CDPSE~·~ISO 27001~·~ISO 22301~·~ISO 9001}\\[0.3ex]
    {\color{white!70!gray}\footnotesize\itshape
        Chercheur — LIMSI — ENSP, Université de Yaoundé~I}\\[0.3ex]
    {\color{white!70!gray}\footnotesize\itshape
        Sr-Eng~·~Sr-Auditor~·~GRC Expert~·~Cryptographe}\\[1.2ex]
    {\color{CouleurAccent}\rule{0.7\linewidth}{1.2pt}}
\end{center}

\vspace{0.6cm}

\begin{center}
    {\color{white!60!gray}\small Année académique 2025--2026}
\end{center}

\afterpage{\nopagecolor}
\end{titlepage}

% ------------------------------------------------------------------------------
% 2. PAGE DE DROITS / VERSO DE LA COUVERTURE
% ------------------------------------------------------------------------------
\clearpage
\thispagestyle{empty}
\vspace*{\fill}

\begin{center}
    \begin{minipage}{0.78\textwidth}
        \footnotesize\color{CouleurBordInfo}
        \begin{center}
            {\bfseries\color{CouleurPrimaire}
                Sécurité Informatique}\\[0.4ex]
            {\itshape Fondamentaux, Pratiques et Gouvernance de la Cybersécurité}
        \end{center}
        \vspace{0.8ex}
        \hrule
        \vspace{0.8ex}
        \textbf{Auteur :} MINKA MI NGUIDJOÏ Thierry Emmanuel\\
        \textbf{Institution :} Université de Yaoundé~I — LIMSI\\
        \textbf{Niveau :} Humanités Numériques — Niveau 3 (Licence 3)\\
        \textbf{Année :} 2025--2026\\[0.5ex]
        Ce document est un support de cours universitaire rédigé à des fins
        pédagogiques. Toute reproduction partielle ou totale à des fins commerciales
        est interdite sans l'autorisation écrite de l'auteur. Les références à des
        produits, organisations ou normes sont citées à titre pédagogique uniquement.\\[0.5ex]
        Les certifications CISA, CISM, CGEIT, CRISC et CDPSE sont des marques
        déposées de l'\textbf{ISACA} (\textit{Information Systems Audit and Control Association}).\\[0.5ex]
        
        \hrule
        \vspace{0.5ex}
        {\centering\itshape
        «~La sécurité n'est jamais un état définitif, mais une exigence permanente de vigilance.~»}
    \end{minipage}
\end{center}

\vspace*{\fill}

% ------------------------------------------------------------------------------
% 3. DÉDICACE (IDENTIQUE AU MODÈLE, ADAPTATION DU TEXTE)
% ------------------------------------------------------------------------------
\clearpage
\thispagestyle{empty}
\vspace*{3cm}
\begin{flushright}
    \begin{minipage}{0.6\textwidth}
        \itshape\large\color{CouleurPrimaire}
        À celles et ceux qui veillent chaque jour\\
        sur la confiance numérique que nous accordons\\
        aux systèmes qui nous gouvernent.\\[1.5ex]
        \normalsize\normalfont\color{CouleurAccent}
        Et à tous les étudiants en sécurité\\
        qui liront ces pages avec la rigueur\\
        qu'ils apporteront demain à leur métier.
    \end{minipage}
\end{flushright}

% ------------------------------------------------------------------------------
% 4. REMERCIEMENTS (IDENTIQUE AU MODÈLE, ADAPTATION DU TEXTE)
% ------------------------------------------------------------------------------
\clearpage
\thispagestyle{empty}
\chapter*{Remerciements}
\addcontentsline{toc}{chapter}{Remerciements}

La conception de ce cours a bénéficié de l'expérience conjointe du chercheur et du praticien. L'auteur exprime sa reconnaissance à \textbf{tous les relecteurs anonymes et non} qui ont contribué à l'enrichissement de cette version. 
Une mention spéciale pour \textbf{le Professeur Flavien Serge Mani Onana} et \textbf{le Professeur Thomas Djotio Ndié}, pour leur encadrement qui nous a permis de confronter plusieurs de nos idées ici développées avec la réalité des exigences académiques et professionnelles.

De même, le Directeur du Contrôle Budgétaire au ministère des finances, \textbf{Madame Tabenyang Augusta épouse Ndjock Arrey}, a été d'un appui significatif dans la rédaction de ce cours, son professionnalisme et sa disponibilité forcent l'admiration.

La matière de ce cours a été nourrie par les missions menées au sein des \textbf{Services du Contrôle Supérieur de l'État}, où la sécurité des systèmes d'information n'est pas une abstraction académique mais une responsabilité institutionnelle quotidienne. De même des prestations comme consultant individuel auprès d'entreprises nationales et internationales, ce qui lui donne une dimension transnationale.

À nos étudiants ...

% ------------------------------------------------------------------------------
% 5. TABLE DES MATIÈRES
% ------------------------------------------------------------------------------
\clearpage
\tableofcontents

% ------------------------------------------------------------------------------
% 6. LISTE DES FIGURES (si nécessaire)
% ------------------------------------------------------------------------------
\clearpage
\listoffigures
\addcontentsline{toc}{chapter}{Liste des figures}

% ------------------------------------------------------------------------------
% 7. LISTE DES TABLEAUX
% ------------------------------------------------------------------------------
\listoftables
\addcontentsline{toc}{chapter}{Liste des tableaux}

% ------------------------------------------------------------------------------
% 8. LISTE DES CODES SOURCE
% ------------------------------------------------------------------------------
\lstlistoflistings
\addcontentsline{toc}{chapter}{Liste des codes source}

% ==============================================================================
% CORPS DU COURS
% ==============================================================================
\mainmatter

% --- Introduction Générale ---
% ==============================================================================
% Fichier : chapitres/00_introduction.tex
% Cours   : Sécurité Informatique — ENSPY, Humanités Numériques Niveau 3
% ==============================================================================

\chapter*{Introduction Générale}
\addcontentsline{toc}{chapter}{Introduction Générale}
\markboth{Introduction Générale}{Introduction Générale}

\begin{boxobjectifs}
\textbf{Objectifs de l'introduction :}
\begin{itemize}
    \item Situer le cours dans le contexte contemporain des systèmes d'information.
    \item Distinguer sécurité informatique et cybersécurité.
    \item Comprendre l'asymétrie des cyberattaques et la notion de cyberguerre.
    \item Connaître les prérequis et la structure globale du cours.
\end{itemize}
\end{boxobjectifs}

\section*{Contexte et Enjeux}

Dans le contexte contemporain des entreprises, les systèmes d'information sont devenus omniprésents, jouant un rôle essentiel dans la facilitation des tâches des employés et l'accélération des processus métier. Cette évolution résulte de l'intégration croissante d'une gamme variée d'équipements informatiques, allant des simples imprimantes aux ordinateurs complexes. Cependant, cette intégration n'est pas sans risques, car elle ouvre la voie à une interconnexion étendue entre ces équipements et les machines externes, offrant ainsi aux individus malveillants des opportunités de vol de données, de perturbation du système ou même d'arrêt complet.

Ce cours adopte une approche holistique en couvrant à la fois les concepts généraux de sécurité informatique ainsi que des éléments spécifiques pertinents pour diverses spécialisations : \textit{data scientist}, \textit{manager des systèmes d'information}, \textit{investigateur numérique}, \textit{community manager}, \textit{spécialiste en intelligence artificielle} ou \textit{expert en cybersécurité}.

\section*{Historique de la Sécurité Informatique}

\subsection*{Les origines (1940--1960)}
La sécurité informatique trouve ses racines dans les années 1940, avec les premiers ordinateurs Colossus et ENIAC. La préoccupation principale était la protection des secrets militaires et gouvernementaux, avec des techniques rudimentaires basées sur des contrôles d'accès physiques et des procédures manuelles.

\subsection*{L'essor des réseaux (1960--1980)}
L'apparition des réseaux informatiques a marqué un tournant majeur. La nécessité de sécuriser les données transmises entre les machines a donné naissance à la sécurité réseau. Les premières technologies — chiffrement, pares-feux — ont été développées durant cette période.

\subsection*{Le premier virus informatique (1988)}
En 1988, le premier virus informatique majeur, connu sous le nom de \og{}Morris\fg{}, a été créé par Robert Tappan Morris à l'Université Cornell. Ce ver a exploité une faille dans Unix pour se propager à travers ARPANET, infectant environ 6~000 ordinateurs et attirant l'attention sur la nécessité de renforcer la sécurité.

\subsection*{L'ère de la cybercriminalité (1980--2000)}
L'expansion d'Internet a ouvert la voie à une nouvelle ère de cybercriminalité : virus, malwares, attaques par déni de service (DoS). L'industrie de la sécurité informatique a connu une croissance exponentielle.

\subsection*{Évolution récente (2000--présent)}
La sécurité informatique couvre aujourd'hui réseaux, systèmes d'information, données et applications. L'intelligence artificielle, l'apprentissage automatique et l'automatisation jouent désormais un rôle crucial dans la lutte contre les cybermenaces.

\section*{Sécurité Informatique vs Cybersécurité}

\begin{boxdef}
\textbf{Cybersécurité :} Se concentre sur la protection des systèmes informatiques et des données dans le \textit{cyberespace} — appareils, réseaux, logiciels, données numériques. L'objectif est de contrer les attaques numériques : virus, malwares, intrusions, piratage.

\textbf{Sécurité Informatique :} Champ plus large. Vise à protéger l'ensemble des systèmes d'information, numériques ou non, contre les menaces physiques (vol, destruction), les erreurs humaines, les défaillances techniques et les catastrophes naturelles. Elle \emph{englobe} la cybersécurité.
\end{boxdef}

\section*{La Cyberguerre et l'Asymétrie des Attaques}

La sécurité informatique concerne aujourd'hui non seulement les entreprises, mais également les États. Leurs systèmes d'information, qui gèrent parfois des infrastructures critiques, peuvent être la cible de cyberattaques orchestrées par des acteurs étatiques ou des groupes criminels organisés.

\begin{boxalert}
\textbf{Asymétrie des cyberattaques :} Contrairement aux guerres physiques où la puissance militaire conventionnelle est déterminante, dans le cyberespace, des pays moins puissants peuvent infliger des dommages significatifs à des nations plus grandes et mieux équipées. Cette asymétrie souligne l'importance de la préparation et de la défense des infrastructures numériques.
\end{boxalert}

\section*{Nécessité et Prérequis du Cours}

La sécurité informatique est devenue une priorité absolue dans un paysage où entreprises, institutions et particuliers sont régulièrement confrontés à des cyberattaques. Ce cours de Licence~3 outille les étudiants pour :
\begin{itemize}
    \item Comprendre les principes fondamentaux de la protection des SI.
    \item Sensibiliser collègues et décideurs aux enjeux de la cybersécurité.
    \item Prévenir les attaques et garantir la résilience des SI.
\end{itemize}

\textbf{Prérequis :} Aucune connaissance préalable en sécurité informatique n'est requise. Une familiarité avec les concepts de base des réseaux informatiques et une compréhension élémentaire des algorithmes de chiffrement et de hachage constituent toutefois des atouts précieux.

\section*{Structure du Cours}

Le cours est subdivisé en neuf (9) livres :

\begin{enumerate}
    \item \textbf{Livre~I :} Définitions, concepts et bases pour le cours
    \item \textbf{Livre~II :} Security and Risk Management
    \item \textbf{Livre~III :} Sécurité des Actifs
    \item \textbf{Livre~IV :} Ingénierie et Architecture de Sécurité
    \item \textbf{Livre~V :} Réseaux et Communications
    \item \textbf{Livre~VI :} Gestion des Identités et des Accès
    \item \textbf{Livre~VII :} Évaluation et Tests de la Sécurité
    \item \textbf{Livre~VIII :} Opérations de Sécurité
    \item \textbf{Livre~IX :} Sécurité des Applications
\end{enumerate}


% -------------------------------------------------------------------
% PARTIE I — Définitions, Concepts et Bases
% -------------------------------------------------------------------
\part{Définitions, Concepts et Bases pour le Cours}

% Chapitre 1 — Définitions et concepts fondamentaux
% ==============================================================================
% Fichier : chapitres/01_definitions_concepts.tex
% Livre I : Définitions, Concepts et Bases pour le Cours
% ==============================================================================

\chapter{Définitions, Concepts et Bases pour le Cours}
\label{chap:definitions}

\begin{boxobjectifs}
\textbf{Objectifs du Livre~I :} Fournir les éléments de compréhension et les notions fondamentales manipulées tout au long du cours. À l'issue de ce livre, l'étudiant maîtrisera le vocabulaire essentiel de la sécurité informatique, de la cybersécurité, de la gestion des risques, des actifs, des réseaux et de l'identité.
\end{boxobjectifs}

% ------------------------------------------------------------------------------
\section{Bases de la Sécurité Informatique}
% ------------------------------------------------------------------------------

\subsection{Confidentialité}
La confidentialité, dans le domaine de la sécurité informatique, se réfère à la protection des informations contre tout accès non autorisé. C'est un principe fondamental qui garantit que seules les personnes ou entités autorisées peuvent accéder à des données sensibles (personnelles, financières, commerciales ou gouvernementales).

\subsection{Disponibilité}
La disponibilité vise à garantir que les services et les systèmes informatiques restent accessibles et fonctionnels à tout moment, selon des critères définis. Le niveau visé varie selon le secteur : dans la santé, une disponibilité de $99{,}999\,\%$ peut être impérative ; dans la finance, des temps d'arrêt planifiés en heures creuses peuvent être tolérés.

\subsection{Intégrité}
L'intégrité des données garantit que celles-ci sont conservées et transmises sans altération non autorisée. Les données doivent rester exactes, cohérentes et fiables tout au long de leur cycle de vie, qu'il s'agisse de stockage ou de transmission.

\subsection{Non-répudiation}
La non-répudiation garantit qu'une fois une action ou une transaction effectuée, les parties impliquées ne peuvent pas la nier a posteriori. L'expéditeur ne peut pas nier avoir envoyé un message, et le destinataire ne peut pas nier l'avoir reçu.

\subsection{Cryptographie}
La cryptographie représente un pilier fondamental de la sécurité informatique. Elle consiste en l'application de techniques mathématiques et algorithmiques pour sécuriser les communications et les informations en les rendant illisibles à toute personne n'ayant pas la clé de déchiffrement.

\subsection{Hachage}
Une fonction de hachage est un algorithme qui prend en entrée des données de n'importe quelle taille et génère une sortie de taille fixe, unique pour chaque entrée (\og{}empreinte\fg{} ou \og{}hash\fg{}). En sécurité, toute modification même minime d'une donnée entraîne un changement complet du hash, permettant de détecter toute altération.

\subsection{Attaques Informatiques}
Une attaque informatique est un événement intentionnel ou non qui vise à compromettre la confidentialité, la disponibilité ou l'intégrité des données et des systèmes. Elle peut prendre la forme d'accès non autorisés, de perturbations de services, de modifications de données ou d'exécutions de code malveillant.

% ------------------------------------------------------------------------------
\section{Bases de la Cybersécurité}
% ------------------------------------------------------------------------------

\subsection{Malware}
Le terme \termecle{malware} désigne un logiciel malveillant conçu pour exécuter des actions nuisibles (copie, suppression ou modification de données, perturbation du système, vol d'informations confidentielles, ouverture de portes dérobées).

\subsection{Vers (Worms)}
Le ver est une forme de malware qui se caractérise par sa capacité à se propager rapidement à travers les réseaux informatiques en se copiant automatiquement dès la détection d'une connexion externe.

\subsection{Ransomware}
Un \termecle{ransomware} chiffre ou verrouille les fichiers ou sous-systèmes de la victime, puis exige le paiement d'une rançon (souvent en cryptomonnaie) pour en restituer l'accès.

\subsection{Ingénierie Sociale}
L'ingénierie sociale exploite les vulnérabilités humaines (peur, cupidité, confiance, soumission à l'autorité) pour tromper les individus et leur faire divulguer des informations confidentielles. Le \termecle{phishing} en est l'exemple le plus répandu.

\subsection{DoS et DDoS}
\begin{itemize}
    \item \textbf{DoS (Denial of Service) :} Attaque visant à rendre un service indisponible en le saturant de requêtes depuis une seule source.
    \item \textbf{DDoS (Distributed DoS) :} Même principe mais depuis des milliers de machines (botnet) simultanément, amplifiant l'impact.
\end{itemize}

\subsection{Botnet}
Un \termecle{botnet} est un réseau de machines infectées (bots) contrôlées à distance par un cybercriminel pour mener des attaques coordonnées (spam, DDoS, minage de cryptomonnaie).

% ------------------------------------------------------------------------------
\section{Bases de la Gestion des Risques}
% ------------------------------------------------------------------------------

\begin{boxdef}
\textbf{Menace :} Tout événement potentiel susceptible de causer un préjudice à un actif.\\[0.3em]
\textbf{Vulnérabilité :} Faiblesse d'un actif ou d'un contrôle pouvant être exploitée par une menace.\\[0.3em]
\textbf{Impact :} Conséquence de la matérialisation d'une menace sur un actif.\\[0.3em]
\textbf{Probabilité :} Vraisemblance qu'une menace se concrétise.\\[0.3em]
\textbf{Risque Informatique :} Combinaison de la probabilité d'occurrence d'une menace et de son impact potentiel sur l'organisation.
\end{boxdef}

La formule de base du risque est :
\[
\text{Risque} = \text{Probabilité} \times \text{Impact}
\]

% ------------------------------------------------------------------------------
\section{Bases de la Gestion des Actifs}
% ------------------------------------------------------------------------------

\begin{boxdef}
\textbf{Actif :} Tout élément ayant de la valeur pour l'organisation et nécessitant une protection.\\[0.3em]
\textbf{Actif Primaire :} Information ou processus métier essentiel (ex. : données clients, processus de paie).\\[0.3em]
\textbf{Actif Secondaire :} Support nécessaire au traitement ou à la protection des actifs primaires (ex. : serveurs, logiciels, personnes).
\end{boxdef}

% ------------------------------------------------------------------------------
\section{Bases en Architecture et Ingénierie de Sécurité}
% ------------------------------------------------------------------------------

\begin{itemize}
    \item \textbf{Architecture de sécurité :} Ensemble cohérent de principes, méthodes et modèles qui guident la conception et le déploiement des mesures de protection d'un SI.
    \item \textbf{Contremesures :} Actions, dispositifs, procédures ou techniques qui réduisent la probabilité ou l'impact d'une menace.
    \item \textbf{Cryptanalyse :} Science visant à analyser et casser les algorithmes de chiffrement, par opposition à la cryptographie.
    \item \textbf{Virtualisation :} Technique consistant à créer des versions virtuelles de ressources matérielles (serveurs, réseaux, stockage) pour optimiser leur utilisation.
\end{itemize}

% ------------------------------------------------------------------------------
\section{Bases des Réseaux et Communications}
% ------------------------------------------------------------------------------

\begin{itemize}
    \item \textbf{Réseaux informatiques :} Ensemble d'équipements interconnectés permettant l'échange de données.
    \item \textbf{Topologie :} Organisation physique ou logique des nœuds d'un réseau (étoile, bus, anneau, maillage).
    \item \textbf{Protocole :} Ensemble de règles définissant le format et l'ordre des messages échangés entre entités réseau.
    \item \textbf{Sécurité réseau :} Ensemble des politiques et mécanismes visant à protéger les données en transit et les équipements réseau.
    \item \textbf{Pare-feu :} Dispositif filtrant le trafic réseau selon des règles définies pour bloquer les accès non autorisés.
    \item \textbf{IDS/IPS :} Systèmes de détection (IDS) et de prévention (IPS) d'intrusion, surveillant le trafic en temps réel.
\end{itemize}

% ------------------------------------------------------------------------------
\section{Bases de la Gestion des Accès et des Identités}
% ------------------------------------------------------------------------------

\begin{itemize}
    \item \textbf{Authentification :} Processus de vérification de l'identité d'un utilisateur ou d'un système.
    \item \textbf{Autorisation :} Processus de détermination des droits et privilèges accordés à une identité authentifiée.
    \item \textbf{Journalisation :} Enregistrement horodaté des événements et actions survenant dans un SI.
    \item \textbf{Contrôle d'accès :} Mécanisme régissant l'accès aux ressources selon le principe du moindre privilège.
\end{itemize}

% ------------------------------------------------------------------------------
\section{Bases de la Sécurité des Opérations}
% ------------------------------------------------------------------------------

\begin{itemize}
    \item \textbf{Incident :} Événement qui compromet ou menace la sécurité d'un SI.
    \item \textbf{Activité :} Tâche ou processus métier réalisé par des personnes ou des systèmes.
    \item \textbf{Opération :} Ensemble d'activités coordonnées pour atteindre un objectif défini.
    \item \textbf{Contrôle :} Mesure technique, organisationnelle ou physique visant à réduire un risque.
\end{itemize}

% ------------------------------------------------------------------------------
\section{Bases du Développement Sécurisé d'Applications}
% ------------------------------------------------------------------------------

\begin{itemize}
    \item \textbf{Besoins fonctionnels :} Fonctions que le logiciel doit remplir (ce qu'il fait).
    \item \textbf{Besoins non fonctionnels :} Contraintes de qualité, de performance, de sécurité et de maintenabilité (comment il le fait).
    \item \textbf{Logiciel :} Programme informatique répondant à un ensemble de besoins définis.
    \item \textbf{Cycle de vie logiciel :} Ensemble des phases de la vie d'un logiciel, de sa conception à son retrait (SDLC).
\end{itemize}

% ------------------------------------------------------------------------------
\section{Évaluation des Acquis --- Livre~I}
% ------------------------------------------------------------------------------

\begin{boxexo}
\textbf{Questions de révision :}
\begin{enumerate}
    \item Quels sont les trois piliers fondamentaux (trilogie CIA) de la sécurité informatique ?
    \item Quelle est la différence entre un virus, un ver et un ransomware ?
    \item Formulez la relation mathématique entre risque, probabilité et impact.
    \item Distinguez actif primaire et actif secondaire en donnant deux exemples de chacun.
    \item Quelle est la différence entre authentification et autorisation ?
    \item En quoi la cryptanalyse est-elle l'inverse de la cryptographie ?
    \item Un employé reçoit un email lui demandant de cliquer sur un lien pour mettre à jour ses identifiants bancaires. De quel type d'attaque s'agit-il ? Justifiez.
    \item Un système informatique connaît une panne pendant 5 minutes par an. Calculez son taux de disponibilité en pourcentage.
\end{enumerate}
\end{boxexo}


% -------------------------------------------------------------------
% PARTIE II — Security and Risk Management
% -------------------------------------------------------------------
\part{Security and Risk Management}

% Chapitre 2 — Gestion des risques de sécurité
% ==============================================================================
% Fichier : chapitres/02_security_risk_management.tex
% Livre II : Security and Risk Management
% ==============================================================================

\chapter{Security and Risk Management}
\label{chap:risk_management}

\begin{boxobjectifs}
\textbf{Objectifs du Livre~II :} Donner les éléments de compréhension liés à la gestion des risques, domaine majeur de la sécurité informatique. À l'issue, l'étudiant saura élaborer une politique de sécurité, mener une analyse des risques, identifier les cadres légaux applicables et distinguer les types d'attaquants.
\end{boxobjectifs}

% ------------------------------------------------------------------------------
\section{Gouvernance de la Sécurité de l'Information}
% ------------------------------------------------------------------------------

\subsection{Élaboration des Politiques de Sécurité}

Une politique de sécurité est un document officiel qui définit les règles et les directives adoptées par une organisation pour gérer et protéger ses actifs informationnels. Elle doit être :

\begin{itemize}
    \item \textbf{Alignée sur les objectifs stratégiques} de l'organisation.
    \item \textbf{Approuvée par la direction} pour lui conférer son autorité.
    \item \textbf{Communiquée} à l'ensemble du personnel.
    \item \textbf{Régulièrement révisée} pour tenir compte de l'évolution des menaces.
\end{itemize}

\subsection{Procédures de Sécurité}

Les procédures de sécurité déclinent opérationnellement les politiques en étapes concrètes. Elles doivent :
\begin{enumerate}
    \item Être développées en concertation avec les parties prenantes.
    \item Faire l'objet de tests et de validations.
    \item Être communiquées et comprises par tout le personnel concerné.
    \item Être soigneusement documentées et archivées.
    \item Être évaluées régulièrement pour identifier les lacunes.
\end{enumerate}

\subsection{Rôles et Responsabilités}

\begin{table}[H]
\centering
\begin{tabular}{p{3.5cm}p{9cm}}
    \toprule
    \textbf{Acteur} & \textbf{Responsabilités} \\
    \midrule
    \textbf{DSI} (Directeur des SI) & Supervision globale du SI, définition de la stratégie sécurité, alignement sur les objectifs métier, allocation des ressources. \\
    \midrule
    \textbf{RSSI} (Resp. Sécurité SI) & Élaboration et mise en œuvre des politiques et procédures, évaluation des risques, supervision des incidents, sensibilisation du personnel. \\
    \midrule
    \textbf{Responsables métier} & Application des politiques dans leur périmètre, remontée des incidents. \\
    \midrule
    \textbf{Auditeurs} & Évaluation indépendante de la conformité et de l'efficacité des contrôles. \\
    \bottomrule
\end{tabular}
\caption{Principaux rôles dans la gouvernance de la sécurité de l'information}
\end{table}

% ------------------------------------------------------------------------------
\section{Contrôle d'Accès}
% ------------------------------------------------------------------------------

Les contrôles d'accès régulent l'accès aux systèmes d'information et se répartissent en deux grandes catégories :

\subsection{Contrôles d'Accès Physique}
\begin{itemize}
    \item Barrières physiques (portiques, tourniquets, sas de sécurité).
    \item Systèmes \og{}\textit{mantrap}\fg{} (accès d'une personne à la fois).
    \item Badges d'accès (RFID, magnétiques).
    \item Surveillance vidéo (CCTV).
    \item Gardiens et rondes de sécurité.
\end{itemize}

\subsection{Contrôles d'Accès Technico-logiques}
\begin{itemize}
    \item Listes de contrôle d'accès (ACL).
    \item Pares-feux filtrant le trafic réseau.
    \item Systèmes IDS/IPS surveillant les activités suspectes.
    \item Authentification multifacteur (MFA).
    \item Chiffrement des données en transit et au repos.
\end{itemize}

\begin{boxinfo}
La combinaison des contrôles physiques et logiques est essentielle pour garantir une sécurité globale. Ils protègent à la fois les infrastructures matérielles et les données numériques.
\end{boxinfo}

% ------------------------------------------------------------------------------
\section{Analyse des Risques}
% ------------------------------------------------------------------------------

\begin{center}
\begin{tikzpicture}[
    node distance=1.5cm,
    box/.style={rectangle, rounded corners, draw=CouleurPrimaire, fill=CouleurPrimaire!10,
                text width=3cm, align=center, minimum height=1cm, font=\small\bfseries}
]
    \node[box] (id)    {Identification\\des risques};
    \node[box] (eval)  [right=of id]   {Évaluation\\des risques};
    \node[box] (trait) [right=of eval] {Traitement\\des risques};
    \node[box] (surv)  [below=of trait] {Surveillance\\et révision};
    \node[box] (comm)  [left=of surv]  {Communication\\des risques};

    \draw[->, thick, CouleurPrimaire] (id) -- (eval);
    \draw[->, thick, CouleurPrimaire] (eval) -- (trait);
    \draw[->, thick, CouleurPrimaire] (trait) -- (surv);
    \draw[->, thick, CouleurPrimaire] (surv) -- (comm);
    \draw[->, thick, CouleurPrimaire] (comm) -- (id);
\end{tikzpicture}
\captionof{figure}{Cycle de gestion des risques}
\end{center}

\subsection{Identification des Risques}

Techniques courantes d'identification :
\begin{itemize}
    \item \textbf{Analyse documentaire :} Examen des politiques, rapports d'incidents et audits passés.
    \item \textbf{Entretiens et ateliers :} Recueil des perceptions des parties prenantes.
    \item \textbf{Analyse de flux :} Cartographie des traitements et des échanges de données.
    \item \textbf{Scans de vulnérabilités :} Outils automatisés identifiant les failles techniques.
\end{itemize}

\subsection{Évaluation des Risques}

\[
\text{Niveau de Risque} = \text{Probabilité d'occurrence} \times \text{Impact potentiel}
\]

\begin{table}[H]
\centering
\small
\begin{tabular}{lccc}
    \toprule
    & \textbf{Impact Faible} & \textbf{Impact Moyen} & \textbf{Impact Élevé} \\
    \midrule
    \textbf{Probabilité Élevée}  & Moyen  & Élevé    & Critique \\
    \textbf{Probabilité Moyenne} & Faible & Moyen    & Élevé    \\
    \textbf{Probabilité Faible}  & Faible & Faible   & Moyen    \\
    \bottomrule
\end{tabular}
\caption{Matrice de criticité des risques}
\end{table}

\subsection{Traitement du Risque}

\begin{itemize}
    \item \textbf{Réduction (Mitigation) :} Mise en place de contrôles pour diminuer la probabilité ou l'impact.
    \item \textbf{Transfert :} Externalisation du risque (assurance, sous-traitance contractuelle).
    \item \textbf{Acceptation :} Décision consciente d'assumer le risque résiduel.
    \item \textbf{Évitement :} Suppression de l'activité ou de l'actif générateur de risque.
\end{itemize}

% ------------------------------------------------------------------------------
\section{Lois et Régulations}
% ------------------------------------------------------------------------------

\subsection{Lois Internationales}
\begin{itemize}
    \item \textbf{RGPD} (Règlement Général sur la Protection des Données, UE 2018) : Protection des données personnelles des citoyens européens.
    \item \textbf{CFAA} (Computer Fraud and Abuse Act, USA) : Cadre pénal américain contre la fraude informatique.
    \item \textbf{Convention de Budapest} (2001) : Premier traité international sur la cybercriminalité.
\end{itemize}

\subsection{Lois Nationales (Cameroun)}
\begin{itemize}
    \item \textbf{Loi n°~2010/012} relative à la cybersécurité et à la cybercriminalité.
    \item Rôle de l'\textbf{ANTIC} (Agence Nationale des Technologies de l'Information et de la Communication) dans la régulation et la supervision de la sécurité des SI.
\end{itemize}

% ------------------------------------------------------------------------------
\section{Sécurité de Tierces Parties}
% ------------------------------------------------------------------------------

\subsection{Responsabilités envers les Tierces Parties}
L'organisation reste responsable de la sécurité des données, même lorsque leur traitement est confié à un tiers (prestataire, fournisseur Cloud, sous-traitant). La due diligence s'impose.

\subsection{Cadre de Gestion des Risques Tiers}
\begin{itemize}
    \item Évaluation de la maturité sécurité du prestataire (questionnaires, audits).
    \item Clauses contractuelles de sécurité (SLA, confidentialité, notification d'incidents).
    \item Surveillance continue des accès du prestataire au SI.
    \item Plan de sortie prévu en cas de défaillance du tiers.
\end{itemize}

% ------------------------------------------------------------------------------
\section{Éthique et Déontologie}
% ------------------------------------------------------------------------------

\begin{boxalert}
\textbf{Code d'éthique du professionnel de sécurité (ISACA/ISC²) :}
\begin{itemize}
    \item Protéger la société, les infrastructures communes et la cybersphère.
    \item Agir honorablement, honnêtement, justement, responsablement et légalement.
    \item Fournir des services diligents et compétents aux clients et aux employeurs.
    \item Promouvoir et protéger la profession.
\end{itemize}
\end{boxalert}

% ------------------------------------------------------------------------------
\section{Types d'Attaquants}
% ------------------------------------------------------------------------------

\begin{table}[H]
\centering
\small
\begin{tabular}{p{3cm}p{4cm}p{5cm}}
    \toprule
    \textbf{Type} & \textbf{Profil} & \textbf{Motivation} \\
    \midrule
    \textbf{Script Kiddie}    & Néophyte utilisant des outils prêts-à-l'emploi & Notoriété, curiosité \\
    \textbf{Hacker Black Hat} & Expert malveillant & Gain financier, espionnage \\
    \textbf{Hacker White Hat} & Expert éthique (pentest autorisé) & Améliorer la sécurité \\
    \textbf{Hacker Grey Hat}  & Entre les deux & Variable \\
    \textbf{Insider Threat}   & Employé malveillant ou négligent & Vengeance, gain, erreur humaine \\
    \textbf{Hacktiviste}      & Groupe idéologique & Protestation politique \\
    \textbf{État-nation}      & Service de renseignement & Espionnage, sabotage stratégique \\
    \bottomrule
\end{tabular}
\caption{Typologie des attaquants}
\end{table}

% ------------------------------------------------------------------------------
\section{Évaluation des Acquis --- Livre~II}
% ------------------------------------------------------------------------------

\begin{boxexo}
\textbf{QCM :}
\begin{enumerate}
    \item La formule $\text{Risque} = \text{Probabilité} \times \text{Impact}$ est utilisée dans quelle étape de la gestion des risques ? \textbf{(a) Identification \quad (b) Évaluation \quad (c) Traitement \quad (d) Surveillance}
    \item Quel est le rôle principal du RSSI ? \textbf{(a) Gérer le budget IT \quad (b) Élaborer et mettre en œuvre la politique sécurité \quad (c) Administrer les serveurs \quad (d) Développer les applications}
    \item Le transfert de risque consiste à : \textbf{(a) Supprimer le risque \quad (b) Le confier à un tiers (assurance) \quad (c) L'ignorer \quad (d) Réduire son impact}
    \item Un attaquant utilisant des outils trouvés sur Internet sans comprendre leur fonctionnement est appelé : \textbf{(a) Hacker Black Hat \quad (b) Insider Threat \quad (c) Script Kiddie \quad (d) Hacktiviste}
\end{enumerate}

\textbf{Travail Pratique :}\\
Réaliser une analyse de risque simplifiée pour un système d'information fictif en utilisant la matrice de criticité. Identifier trois risques, les évaluer et proposer un traitement pour chacun.
\end{boxexo}


% -------------------------------------------------------------------
% PARTIE III — Sécurité des Actifs
% -------------------------------------------------------------------
\part{Sécurité des Actifs}

% Chapitre 3 — Classification et protection des actifs
% ==============================================================================
% Fichier : chapitres/03_securite_actifs.tex
% Livre III : Sécurité des Actifs
% ==============================================================================

\chapter{Sécurité des Actifs}
\label{chap:securite_actifs}

\begin{boxobjectifs}
\textbf{Objectifs du Livre~III :} Donner les éléments de compréhension liés aux actifs informationnels, leur importance et leur protection. À l'issue, l'étudiant saura classifier les données, appliquer les modèles de contrôle d'accès appropriés, gérer la rémanence et se référer aux standards internationaux de contrôle.
\end{boxobjectifs}

% ------------------------------------------------------------------------------
\section{Classification des Données}
% ------------------------------------------------------------------------------

\subsection{Niveaux de Sensibilité}

La classification des données permet d'appliquer des mesures de protection proportionnées à leur valeur et à leur sensibilité.

\begin{table}[H]
\centering
\begin{tabular}{p{3cm}p{5cm}p{4cm}}
    \toprule
    \textbf{Niveau} & \textbf{Description} & \textbf{Exemples} \\
    \midrule
    \textbf{Public}        & Aucune restriction de diffusion. & Site web institutionnel, communiqués de presse. \\
    \textbf{Interne}       & Usage réservé aux employés. & Procédures internes, organigrammes. \\
    \textbf{Confidentiel}  & Accès restreint aux personnes autorisées. & Contrats clients, données RH. \\
    \textbf{Secret/Très Sensible} & Accès ultra-restreint, impact majeur en cas de divulgation. & Données financières stratégiques, secrets industriels. \\
    \bottomrule
\end{tabular}
\caption{Niveaux de classification des données}
\end{table}

\subsection{Méthodes de Classification}
\begin{itemize}
    \item \textbf{Classification manuelle :} Un responsable désigne le niveau lors de la création.
    \item \textbf{Classification automatique :} Outils DLP (Data Loss Prevention) analysant le contenu et attribuant automatiquement un niveau.
    \item \textbf{Classification basée sur les rôles :} Le niveau dépend du rôle de l'utilisateur créateur.
\end{itemize}

% ------------------------------------------------------------------------------
\section{Protection des Données}
% ------------------------------------------------------------------------------

\subsection{Modèles de Contrôles d'Accès}

\begin{table}[H]
\centering
\small
\begin{tabular}{p{2.5cm}p{4.5cm}p{4.5cm}}
    \toprule
    \textbf{Modèle} & \textbf{Principe} & \textbf{Usage typique} \\
    \midrule
    \textbf{MAC} (Mandatory) & Niveau de sécurité attribué aux sujets et objets par l'administrateur système. L'accès dépend de la comparaison des étiquettes. & Systèmes gouvernementaux, militaires. \\
    \midrule
    \textbf{DAC} (Discretionary) & Le propriétaire de la ressource définit qui peut y accéder. & Systèmes de fichiers courants (Windows NTFS, Linux). \\
    \midrule
    \textbf{RBAC} (Role-Based) & Les droits sont attribués à des rôles ; les utilisateurs héritent des droits de leur rôle. & Applications d'entreprise (ERP, CRM). \\
    \midrule
    \textbf{ABAC} (Attribute-Based) & Accès basé sur les attributs du sujet, de l'objet et de l'environnement. & Cloud, IoT, systèmes distribués complexes. \\
    \bottomrule
\end{tabular}
\caption{Modèles de contrôle d'accès}
\end{table}

\subsection{Mesures Technico-logiques}
\begin{itemize}
    \item \textbf{Pare-feux :} Filtrage du trafic réseau.
    \item \textbf{Antivirus / EDR :} Détection et neutralisation des malwares.
    \item \textbf{Chiffrement :} Protection des données en transit (TLS) et au repos (AES).
    \item \textbf{DLP (Data Loss Prevention) :} Surveillance et blocage des fuites de données.
\end{itemize}

% ------------------------------------------------------------------------------
\section{Appartenance des Données (Data Ownership)}
% ------------------------------------------------------------------------------

\subsection{Définition et Importance}
L'appartenance des données désigne la propriété et la responsabilité de celles-ci au sein d'une organisation. Identifier clairement le propriétaire d'une donnée permet de :
\begin{itemize}
    \item Définir les rôles et responsabilités en matière de gestion.
    \item Établir des droits d'accès appropriés.
    \item Faciliter la traçabilité.
    \item Assurer la conformité réglementaire.
\end{itemize}

\subsection{Méthodes pour Déterminer l'Appartenance}
\begin{itemize}
    \item Identification du propriétaire (créateur ou responsable métier de la donnée).
    \item Analyse des flux de données au sein de l'organisation.
    \item Exploitation des métadonnées (auteur, date, département).
\end{itemize}

\subsection{Bonnes Pratiques}
\begin{itemize}
    \item Sensibiliser les employés à leurs responsabilités.
    \item Définir des politiques de classification et de gestion des accès claires.
    \item Réaliser des audits réguliers de la gestion des accès.
\end{itemize}

% ------------------------------------------------------------------------------
\section{Mémoire, Rémanence et Destruction}
% ------------------------------------------------------------------------------

\begin{boxdef}
\textbf{Rémanence des données :} Persistance de données dans un support de stockage, même après leur suppression apparente. Ce phénomène peut conduire à des fuites d'informations sensibles si les supports ne sont pas correctement effacés avant réaffectation ou destruction.
\end{boxdef}

\subsection{Risques associés à la Rémanence}
\begin{itemize}
    \item Vol de données lors de la revente ou du recyclage de matériel.
    \item Réutilisation inappropriée de supports sans effacement sécurisé.
    \item Non-conformité réglementaire (conservation de données au-delà de la durée légale).
\end{itemize}

\subsection{Méthodes d'Effacement Sécurisé}

\begin{table}[H]
\centering
\begin{tabular}{p{3cm}p{9cm}}
    \toprule
    \textbf{Méthode} & \textbf{Description} \\
    \midrule
    \textbf{Écrasement}   & Réécriture multiple sur les secteurs (normes DoD~5220.22-M : 3 passes minimum). \\
    \textbf{Dégaussement} & Exposition à un champ magnétique puissant détruisant les données magnétiques. \\
    \textbf{Destruction physique} & Déchiquetage, broyage ou incinération du support. Méthode la plus sûre. \\
    \textbf{Chiffrement + effacement} & Chiffrement préalable des données puis suppression de la clé. \\
    \bottomrule
\end{tabular}
\caption{Méthodes d'effacement sécurisé des données}
\end{table}

% ------------------------------------------------------------------------------
\section{Standards et Référentiels de Contrôle}
% ------------------------------------------------------------------------------

\begin{table}[H]
\centering
\small
\begin{tabular}{p{3.5cm}p{8.5cm}}
    \toprule
    \textbf{Standard} & \textbf{Description} \\
    \midrule
    \textbf{ISO/IEC~27001} & Système de Management de la Sécurité de l'Information (SMSI). Certification internationale de référence. \\
    \textbf{NIST SP~800-53} & Contrôles de sécurité pour les systèmes d'information fédéraux américains (18 familles de contrôles). \\
    \textbf{COBIT~2019}     & Cadre de gouvernance et de gestion des TI développé par l'ISACA (5 domaines). \\
    \textbf{CIS Controls}   & 18 contrôles de sécurité essentiels prioritaires, du Center for Internet Security. \\
    \textbf{PCI~DSS}        & Norme de sécurité des données pour l'industrie des cartes de paiement (12 exigences). \\
    \bottomrule
\end{tabular}
\caption{Principaux standards et référentiels de contrôle}
\end{table}

% ------------------------------------------------------------------------------
\section{Évaluation des Acquis --- Livre~III}
% ------------------------------------------------------------------------------

\begin{boxexo}
\textbf{QCM :}
\begin{enumerate}
    \item Quel modèle de contrôle d'accès est le plus adapté à un système militaire ? \textbf{(a) DAC \quad (b) RBAC \quad (c) MAC \quad (d) ABAC}
    \item La rémanence des données désigne : \textbf{(a) La durée de vie d'un disque dur \quad (b) La persistance de données après leur suppression apparente \quad (c) La sauvegarde automatique des données \quad (d) La compression des données}
    \item ISO/IEC 27001 porte sur : \textbf{(a) Les contrôles financiers \quad (b) Le système de management de la sécurité de l'information \quad (c) La gestion de projet \quad (d) La virtualisation}
    \item Quelle méthode d'effacement est la plus sûre pour un disque dur contenant des données classifiées Secret ? \textbf{(a) Formatage simple \quad (b) Suppression des fichiers \quad (c) Destruction physique \quad (d) Chiffrement ultérieur}
\end{enumerate}

\textbf{Travail Pratique :}\\
Proposer une politique de classification des données pour une université, en identifiant au moins cinq types de données traitées, leurs niveaux de sensibilité et les mesures de protection à appliquer.
\end{boxexo}


% -------------------------------------------------------------------
% PARTIE IV — Ingénierie et Architecture de Sécurité
% -------------------------------------------------------------------
\part{Ingénierie et Architecture de Sécurité}

% Chapitre 4 — Principes d'architecture sécurisée
% ==============================================================================
% Fichier : chapitres/04_ingenierie_architecture.tex
% Livre IV : Ingénierie et Architecture de Sécurité
% ==============================================================================

\chapter{Ingénierie et Architecture de Sécurité}
\label{chap:ingenierie_archi}

\begin{boxobjectifs}
\textbf{Objectifs du Livre~IV :} Donner les éléments liés aux modèles de sécurité, aux fondements de la cryptographie (symétrique, asymétrique, hachage), à la sécurité physique et à la virtualisation. À l'issue, l'étudiant sera capable d'implémenter les chiffrements classiques et de concevoir une architecture de sécurité défense en profondeur.
\end{boxobjectifs}

% ------------------------------------------------------------------------------
\section{Modèles de Sécurité}
% ------------------------------------------------------------------------------

\subsection{Modèle de Bell-LaPadula (Confidentialité)}

\begin{boxdef}
Modèle orienté \textbf{confidentialité}, utilisé dans les systèmes militaires. Repose sur deux règles :
\begin{itemize}
    \item \textbf{No Read Up (NRU) :} Un sujet ne peut pas lire une ressource de niveau de classification supérieur au sien.
    \item \textbf{No Write Down (NWD) :} Un sujet ne peut pas écrire dans une ressource de niveau inférieur au sien (pour éviter la fuite d'informations classifiées).
\end{itemize}
\end{boxdef}

\subsection{Modèle de Biba (Intégrité)}
\begin{boxdef}
Modèle orienté \textbf{intégrité}, inverse de Bell-LaPadula :
\begin{itemize}
    \item \textbf{No Write Up :} Un sujet ne peut pas écrire dans une ressource de niveau d'intégrité supérieur.
    \item \textbf{No Read Down :} Un sujet ne peut pas lire une ressource de niveau d'intégrité inférieur.
\end{itemize}
\end{boxdef}

\subsection{Chinese Wall (Brewer-Nash)}
Modèle dynamique évitant les \textbf{conflits d'intérêts}. Un consultant ayant accédé aux données d'une entreprise A ne peut plus accéder aux données d'une entreprise concurrente B.

% ------------------------------------------------------------------------------
\section{Fondamentaux de la Conception d'Architecture de Sécurité}
% ------------------------------------------------------------------------------

\subsection{Principes de Base}

\begin{itemize}
    \item \textbf{Défense en profondeur :} Multiplier les couches de sécurité indépendantes. Si une couche est compromise, les suivantes assurent la protection.
    \item \textbf{Moindre privilège :} N'accorder que les droits strictement nécessaires à l'accomplissement d'une tâche.
    \item \textbf{Séparation des fonctions :} Aucun individu ne doit avoir le contrôle complet sur un processus critique (principe des «~quatre yeux~»).
    \item \textbf{Sécurité par défaut :} Les systèmes doivent être sécurisés dans leur configuration initiale.
    \item \textbf{Fail Secure :} En cas de panne, le système doit basculer dans un état sécurisé (accès refusé plutôt qu'accordé).
\end{itemize}

\subsection{Segmentation Réseau et Défense en Profondeur}

\begin{center}
\begin{tikzpicture}[scale=0.85]
  \draw[thick, CouleurPrimaire, fill=CouleurPrimaire!5, rounded corners] (-5,-3) rectangle (5,3);
  \node[CouleurPrimaire, font=\small\bfseries] at (0,2.7) {Zone Externe (Internet)};

  \draw[thick, CouleurAccent, fill=CouleurAccent!5, rounded corners] (-3.5,-2) rectangle (3.5,2);
  \node[CouleurAccent, font=\small\bfseries] at (0,1.7) {DMZ (Serveurs Web, Mail)};

  \draw[thick, NavyBlue, fill=NavyBlue!5, rounded corners] (-2,-1) rectangle (2,1);
  \node[NavyBlue, font=\small\bfseries] at (0,0.7) {LAN Interne};
  \node[font=\tiny] at (0,0.1) {Serveurs applicatifs, BDD};
  \node[font=\tiny] at (0,-0.3) {Postes de travail};

  % Pares-feux
  \node[font=\tiny\bfseries, CouleurPrimaire] at (4,-2.5) {Pare-feu 1};
  \node[font=\tiny\bfseries, CouleurAccent]  at (3,-1.5) {Pare-feu 2};
\end{tikzpicture}
\captionof{figure}{Architecture Défense en Profondeur avec DMZ}
\end{center}

% ------------------------------------------------------------------------------
\section{Architecture Matérielle Sécurisée}
% ------------------------------------------------------------------------------

\subsection{Vulnérabilités et Menaces Matérielles}
\begin{itemize}
    \item Défauts de conception ou de fabrication (ex. : Spectre, Meltdown).
    \item Configurations incorrectes (BIOS/UEFI non protégés, ports USB ouverts).
    \item Menaces physiques : vol, sabotage, espionnage industriel.
\end{itemize}

\subsection{Mesures de Protection Matérielle}
\begin{itemize}
    \item \textbf{TPM (Trusted Platform Module) :} Puce sécurisée pour le stockage de clés de chiffrement et la vérification de l'intégrité du démarrage.
    \item \textbf{Contrôle d'accès physique :} Salles serveurs verrouillées, racks sécurisés.
    \item \textbf{Surveillance de l'intégrité :} Journaux événements, IDS matériels.
    \item \textbf{Mises à jour de micrologiciel :} Application régulière des patchs BIOS/UEFI.
\end{itemize}

% ------------------------------------------------------------------------------
\section{Systèmes d'Exploitation et Architecture Logicielle}
% ------------------------------------------------------------------------------

\subsection{Vulnérabilités des Systèmes d'Exploitation}
\begin{itemize}
    \item Erreurs de programmation (buffer overflow, use-after-free).
    \item Droits d'accès inappropriés (comptes sur-privilégiés).
    \item Services non utilisés exposant des surfaces d'attaque inutiles.
    \item Mises à jour non appliquées (exploitation des CVE connus).
\end{itemize}

\subsection{Principes d'Architecture Logicielle Sécurisée}
\begin{itemize}
    \item \textbf{Modularité :} Composants indépendants et substituables.
    \item \textbf{Séparation des préoccupations :} Fonctions de sécurité isolées du code métier.
    \item \textbf{Abstraction :} Masquage des détails d'implémentation.
    \item \textbf{Évolutivité :} Capacité à gérer une charge croissante sans dégrader la sécurité.
\end{itemize}

% ------------------------------------------------------------------------------
\section{Virtualisation}
% ------------------------------------------------------------------------------

La virtualisation utilise un \termecle{hyperviseur} pour créer et gérer des machines virtuelles (VM). Elle introduit de nouvelles vulnérabilités :
\begin{itemize}
    \item \textbf{Attaques contre l'hyperviseur (VM Escape) :} Compromission de toutes les VM hébergées.
    \item \textbf{Attaques inter-VM (VM Hopping) :} Une VM compromise attaque les voisines.
    \item \textbf{Fuites de données mémoire} entre VM partageant le même hôte physique.
\end{itemize}

\textbf{Contre-mesures :} Sécurisation de l'hyperviseur, isolation des réseaux virtuels (VLAN), application stricte des politiques de groupes de sécurité.

% ------------------------------------------------------------------------------
\section{Introduction à la Cryptographie}
% ------------------------------------------------------------------------------

\subsection{Algorithmes Symétriques}

La même clé est utilisée pour chiffrer et déchiffrer. Alice et Bob doivent partager la clé secrète à l'avance.

\subsubsection{Chiffrement de César}
Substitution mono-alphabétique par décalage d'un entier $k$ ($0 \leq k \leq 25$) dans l'alphabet.

\[
C_i = (P_i + k) \mod 26 \qquad \text{(Chiffrement)}
\]
\[
P_i = (C_i - k) \mod 26 \qquad \text{(Déchiffrement)}
\]

\begin{boxlizbiz}
\textbf{Exemple :} $k=3$, message \texttt{BONJOUR} $\rightarrow$ \texttt{ERQMRXU}\\
\textit{Faiblesse :} 25 clés seulement, cassé par analyse de fréquence.
\end{boxlizbiz}

\subsubsection{Chiffrement de Vigenère}
Substitution poly-alphabétique utilisant une clé de longueur $n$. La lettre en position $i$ est décalée de la valeur numérique de la lettre de clé correspondante.

\[
C_i = (P_i + K_{i \mod n}) \mod 26
\]

\begin{boxlizbiz}
\textbf{Exemple :} Message \texttt{BONJOUR}, Clé \texttt{CLE} :
\begin{verbatim}
B O N J O U R
C L E C L E C
D Z R L Z Y T
\end{verbatim}
\textit{Faiblesse :} Cassé par le test de Kasiski si la clé est courte et répétée.
\end{boxlizbiz}

\subsubsection{Chiffrement de Hill}
Chiffrement par blocs utilisant l'algèbre matricielle sur $\mathbb{Z}_{26}$ :
\[
\mathbf{C} = \mathbf{K} \cdot \mathbf{P} \mod 26
\]
où $\mathbf{K}$ est la matrice-clé (inversible sur $\mathbb{Z}_{26}$) et $\mathbf{P}$ le vecteur plaintext.

\subsubsection{Algorithmes Modernes Symétriques}
\begin{itemize}
    \item \textbf{DES} (Data Encryption Standard, 1977) : 56 bits, obsolète.
    \item \textbf{3DES} : Triple DES, plus sûr mais lent.
    \item \textbf{AES} (Advanced Encryption Standard, 2001) : Standard actuel — 128, 192 ou 256 bits.
\end{itemize}

\subsection{Algorithmes Asymétriques}
Deux clés différentes : clé publique (chiffrement, vérification de signature) et clé privée (déchiffrement, signature). La clé publique peut être librement partagée.

\begin{itemize}
    \item \textbf{RSA :} Basé sur la difficulté de factorisation de grands entiers. Clés de 2048 ou 4096 bits.
    \item \textbf{ECDSA :} Basé sur les courbes elliptiques. Efficace pour les signatures.
    \item \textbf{Diffie-Hellman :} Protocole d'échange de clés sur un canal non sécurisé.
\end{itemize}

\begin{boxinfo}
\textbf{Utilisation hybride :} En pratique, on utilise la cryptographie asymétrique pour échanger une clé symétrique de session, puis AES pour chiffrer les données (TLS/SSL).
\end{boxinfo}

\subsection{Fonctions de Hachage}
Transforme des données de taille quelconque en empreinte de taille fixe. Propriétés : déterminisme, résistance aux préimages, résistance aux collisions, effet avalanche.

\begin{table}[H]
\centering
\begin{tabular}{llll}
    \toprule
    \textbf{Algorithme} & \textbf{Taille du hash} & \textbf{Statut} \\
    \midrule
    MD5    & 128 bits & Obsolète (collisions connues) \\
    SHA-1  & 160 bits & Déprécié \\
    SHA-256 & 256 bits & Recommandé \\
    SHA-3  & 256/512 bits & Standard actuel \\
    \bottomrule
\end{tabular}
\caption{Fonctions de hachage courantes}
\end{table}

\subsection{Cryptanalyse}
Science visant à analyser et casser les systèmes cryptographiques. Principales attaques :
\begin{itemize}
    \item \textbf{Force brute :} Essai exhaustif de toutes les clés possibles.
    \item \textbf{Analyse de fréquence :} Exploitation des redondances statistiques de la langue.
    \item \textbf{Attaque à texte clair connu :} Exploitation de paires (texte clair, texte chiffré) connues.
    \item \textbf{Attaque par collision :} Recherche de deux entrées distinctes produisant le même hash.
\end{itemize}

% ------------------------------------------------------------------------------
\section{Sécurité Physique}
% ------------------------------------------------------------------------------
\begin{itemize}
    \item Contrôle d'accès aux salles serveurs (badge, biométrie, sas).
    \item Surveillance vidéo et enregistrement des entrées.
    \item Protection contre les incendies (détection, suppression gaz FM200).
    \item Alimentation sans coupure (UPS) et générateurs.
    \item Protection contre les inondations, la chaleur et les interférences électromagnétiques.
\end{itemize}

% ------------------------------------------------------------------------------
\section{Évaluation des Acquis --- Livre~IV}
% ------------------------------------------------------------------------------

\begin{boxexo}
\textbf{Exercices :}
\begin{enumerate}
    \item Chiffrez le message \texttt{ENSPY} avec le chiffrement de César, $k=5$.
    \item Avec la clé \texttt{SEC}, chiffrez \texttt{AUDIT} avec Vigenère.
    \item Quelle est la différence fondamentale entre un algorithme symétrique et un algorithme asymétrique ?
    \item Expliquez le principe de défense en profondeur avec un schéma réseau commenté.
    \item Calculez le hash SHA-256 de la chaîne \texttt{Securite2024} (à faire en TP avec Python : \texttt{hashlib.sha256}).
\end{enumerate}

\textbf{Travail Pratique :}\\
Concevoir et développer un site web sécurisé implémentant le modèle Chinese Wall pour une entreprise de traitement de données. Le site doit : (1) authentifier les utilisateurs avec un hash sécurisé (bcrypt), (2) vérifier l'intégrité des données (SHA-256), (3) appliquer les restrictions Chinese Wall selon les projets affectés.
\end{boxexo}


% -------------------------------------------------------------------
% PARTIE V — Réseaux et Communications
% -------------------------------------------------------------------
\part{Réseaux et Communications}

% Chapitre 5 — Sécurité des réseaux et communications
% ==============================================================================
% Fichier : chapitres/05_reseaux_communications.tex
% Livre V : Réseaux et Communications
% ==============================================================================

\chapter{Réseaux et Communications}
\label{chap:reseaux}

\begin{boxobjectifs}
\textbf{Objectifs du Livre~V :} Rappeler les fondamentaux des modèles et protocoles réseaux. Présenter les méthodes et mécanismes de sécurisation des réseaux et des communications (VPN, certificats, PKI). Identifier les principales attaques réseau et leurs contre-mesures.
\end{boxobjectifs}

% ------------------------------------------------------------------------------
\section{Les Modèles Réseaux}
% ------------------------------------------------------------------------------

\subsection{Modèle OSI (7 couches)}

\begin{table}[H]
\centering
\small
\begin{tabular}{clp{4cm}p{4cm}}
    \toprule
    \textbf{N°} & \textbf{Couche} & \textbf{Rôle} & \textbf{Protocoles / Équipements} \\
    \midrule
    7 & Application    & Services aux applications & HTTP, FTP, DNS, SMTP, SSH \\
    6 & Présentation   & Format, chiffrement, compression & SSL/TLS, JPEG, ASCII \\
    5 & Session        & Établissement, maintien, fin de sessions & NetBIOS, RPC \\
    4 & Transport      & Fiabilité, contrôle de flux & TCP, UDP \\
    3 & Réseau         & Adressage logique, routage & IP, ICMP, OSPF \\
    2 & Liaison        & Adressage physique, détection d'erreurs & Ethernet, Wi-Fi (802.11), ARP \\
    1 & Physique       & Transmission des bits & Câbles, fibre optique, hubs \\
    \bottomrule
\end{tabular}
\caption{Modèle OSI}
\end{table}

\subsection{Modèle TCP/IP (4 couches)}

\begin{table}[H]
\centering
\begin{tabular}{clp{7cm}}
    \toprule
    \textbf{N°} & \textbf{Couche} & \textbf{Équivalence OSI} \\
    \midrule
    4 & Application & Couches 5, 6, 7 \\
    3 & Transport   & Couche 4 \\
    2 & Internet    & Couche 3 \\
    1 & Accès réseau & Couches 1 et 2 \\
    \bottomrule
\end{tabular}
\caption{Modèle TCP/IP}
\end{table}

% ------------------------------------------------------------------------------
\section{Protocoles des Couches Basses}
% ------------------------------------------------------------------------------

\begin{itemize}
    \item \textbf{ARP (Address Resolution Protocol) :} Résolution d'une adresse IP en adresse MAC. Vulnérable à l'ARP Spoofing.
    \item \textbf{ICMP :} Protocole de diagnostic réseau (\texttt{ping}). Peut être utilisé pour des attaques de reconnaissance.
    \item \textbf{TCP (Transmission Control Protocol) :} Connexion fiable, orientée flux. Handshake SYN-SYN/ACK-ACK. Vulnérable au SYN Flood.
    \item \textbf{UDP (User Datagram Protocol) :} Sans connexion, non fiable, rapide. Utilisé pour DNS, streaming, VoIP.
\end{itemize}

% ------------------------------------------------------------------------------
\section{Protocoles Applicatifs Fondamentaux}
% ------------------------------------------------------------------------------

\begin{table}[H]
\centering
\small
\begin{tabular}{p{2.5cm}p{1.5cm}p{7cm}}
    \toprule
    \textbf{Protocole} & \textbf{Port} & \textbf{Description et risques} \\
    \midrule
    HTTP  & 80  & Web non sécurisé. Trafic en clair, vulnérable à l'interception. \\
    HTTPS & 443 & HTTP sur TLS. Chiffrement du trafic web. Standard actuel. \\
    FTP   & 21  & Transfert de fichiers en clair. Remplacé par SFTP/FTPS. \\
    SSH   & 22  & Shell sécurisé chiffré. Remplace Telnet (23, en clair). \\
    DNS   & 53  & Résolution de noms. Vulnérable au DNS Spoofing et Cache Poisoning. \\
    SMTP  & 25  & Envoi d'emails. Vecteur de phishing et spam. \\
    LDAP  & 389 & Annuaire d'authentification. LDAPS (636) pour la version sécurisée. \\
    \bottomrule
\end{tabular}
\caption{Protocoles applicatifs courants et risques associés}
\end{table}

% ------------------------------------------------------------------------------
\section{Technologies LAN et WAN}
% ------------------------------------------------------------------------------

\begin{itemize}
    \item \textbf{VLAN (Virtual LAN) :} Segmentation logique d'un réseau physique. Permet d'isoler des zones (ex : VLAN Finance séparé du VLAN Invités).
    \item \textbf{DMZ (Demilitarized Zone) :} Zone tampon entre Internet et le LAN interne, hébergeant les serveurs exposés (web, mail).
    \item \textbf{WAN / MPLS :} Interconnexion de sites distants de manière privée et sécurisée.
\end{itemize}

% ------------------------------------------------------------------------------
\section{Périphériques Réseau}
% ------------------------------------------------------------------------------

\begin{table}[H]
\centering
\small
\begin{tabular}{p{3cm}p{4cm}p{5cm}}
    \toprule
    \textbf{Équipement} & \textbf{Rôle} & \textbf{Aspect sécurité} \\
    \midrule
    \textbf{Routeur}    & Routage entre réseaux (Couche 3) & ACL, filtrage de paquets \\
    \textbf{Commutateur (Switch)} & Commutation dans un LAN (Couche 2) & Port security, désactivation VTP, 802.1X \\
    \textbf{Pare-feu}   & Filtrage du trafic selon des règles & Stateful inspection, NAT, IPS intégré \\
    \textbf{IDS/IPS}    & Détection/Prévention d'intrusion & Signatures, anomalies comportementales \\
    \textbf{Proxy}      & Intermédiaire pour les requêtes web & Filtrage URL, contrôle du contenu \\
    \bottomrule
\end{tabular}
\caption{Équipements réseau et sécurité associée}
\end{table}

% ------------------------------------------------------------------------------
\section{Communications Sécurisées}
% ------------------------------------------------------------------------------

\subsection{Protocoles de Sécurité}

\begin{itemize}
    \item \textbf{TLS/SSL :} Chiffrement des communications web (HTTPS). TLS~1.3 est le standard actuel.
    \item \textbf{IPSec :} Sécurisation au niveau IP (chiffrement + authentification). Utilisé dans les VPN.
    \item \textbf{SSH :} Accès distant sécurisé, tunneling.
    \item \textbf{SFTP/FTPS :} Transfert de fichiers sécurisé.
\end{itemize}

\subsection{Certificats Numériques et PKI}

\begin{boxdef}
\textbf{Certificat numérique :} Document électronique signé par une Autorité de Certification (AC) qui lie une clé publique à une identité. Format standard : X.509.

\textbf{PKI (Public Key Infrastructure) :} Ensemble des politiques, procédures, rôles et technologies permettant de gérer les certificats numériques. Composants : AC Racine, AC Intermédiaire, Annuaire LDAP, CRL (Certificate Revocation List).
\end{boxdef}

\subsection{VPN (Virtual Private Network)}
Le VPN crée un tunnel chiffré sur un réseau non sécurisé (Internet), permettant une connexion sécurisée à distance. Deux types principaux :
\begin{itemize}
    \item \textbf{VPN Site-à-Site :} Interconnexion permanente de deux sites distants.
    \item \textbf{VPN Accès distant (SSL/TLS) :} Connexion ponctuelle d'un utilisateur nomade au SI de l'entreprise.
\end{itemize}

% ------------------------------------------------------------------------------
\section{Attaques Réseau}
% ------------------------------------------------------------------------------

\begin{table}[H]
\centering
\small
\begin{tabular}{p{3cm}p{5cm}p{4cm}}
    \toprule
    \textbf{Attaque} & \textbf{Description} & \textbf{Contre-mesure} \\
    \midrule
    \textbf{ARP Spoofing} & Empoisonnement du cache ARP pour intercepter le trafic (MITM). & ARP inspection dynamique (DAI), segmentation VLAN. \\
    \textbf{DNS Spoofing} & Réponses DNS falsifiées redirigeant vers des serveurs malveillants. & DNSSEC, validation des réponses. \\
    \textbf{SYN Flood} & Saturation de la table de connexions TCP par des SYN sans ACK final. & SYN cookies, limitation de connexions. \\
    \textbf{Man-in-the-Middle} & Interception des communications entre deux parties. & Chiffrement TLS, vérification des certificats. \\
    \textbf{Sniffer/Wireshark} & Capture passive du trafic réseau en clair. & Chiffrement (TLS/SSH), switch au lieu de hub. \\
    \textbf{Port Scanning} & Cartographie des ports ouverts (Nmap). & Pare-feu, masquage des bannières de services. \\
    \bottomrule
\end{tabular}
\caption{Principales attaques réseau et contre-mesures}
\end{table}

% ------------------------------------------------------------------------------
\section{Évaluation des Acquis --- Livre~V}
% ------------------------------------------------------------------------------

\begin{boxexo}
\textbf{QCM :}
\begin{enumerate}
    \item À quelle couche OSI le protocole IP opère-t-il ? \textbf{(a) 2 \quad (b) 3 \quad (c) 4 \quad (d) 7}
    \item Lequel des protocoles suivants chiffre les communications ? \textbf{(a) HTTP \quad (b) FTP \quad (c) Telnet \quad (d) SSH}
    \item Un ensemble d'équipements ou de services exposés à Internet (serveurs web, mail), séparés du réseau interne par un pare-feu, s'appelle : \textbf{(a) VLAN \quad (b) DMZ \quad (c) VPN \quad (d) NAT}
    \item L'attaque consistant à répondre à des requêtes ARP avec de fausses adresses MAC est : \textbf{(a) SYN Flood \quad (b) DNS Spoofing \quad (c) ARP Spoofing \quad (d) Port Scanning}
\end{enumerate}

\textbf{Travail Pratique :}\\
Sur Cisco Packet Tracer, monter le réseau de l'entreprise VISION composé de deux sites distants reliés par un VPN. Configurer une DMZ hébergeant les serveurs web et mail. Tester la connectivité et l'isolation de la DMZ vis-à-vis du LAN interne.
\end{boxexo}


% -------------------------------------------------------------------
% PARTIE VI — Gestion des Identités et des Accès
% -------------------------------------------------------------------
\part{Gestion des Identités et des Accès}

% Chapitre 6 — IAM et contrôle d'accès
% ==============================================================================
% Fichier : chapitres/06_gestion_identites_acces.tex
% Livre VI : Gestion des Identités et des Accès
% ==============================================================================

\chapter{Gestion des Identités et des Accès}
\label{chap:iam}

\begin{boxobjectifs}
\textbf{Objectifs du Livre~VI :} Maîtriser le cadre IAAA, les méthodes d'authentification, les modèles de contrôle d'accès et l'implémentation des contrôles dans les organisations.
\end{boxobjectifs}

% ------------------------------------------------------------------------------
\section{Le Cadre IAAA}
% ------------------------------------------------------------------------------

\begin{boxdef}
\begin{itemize}
    \item \textbf{Identification :} Une entité prétend posséder une identité (ex. : saisie d'un nom d'utilisateur).
    \item \textbf{Authentification :} Vérification de cette identité (ex. : mot de passe, OTP, biométrie).
    \item \textbf{Autorisation :} Détermination des droits et privilèges accordés à l'identité authentifiée.
    \item \textbf{Accountability (Imputabilité) :} Traçabilité de toutes les actions effectuées par l'identité (journalisation, audit).
\end{itemize}
\end{boxdef}

% ------------------------------------------------------------------------------
\section{Méthodes d'Authentification}
% ------------------------------------------------------------------------------

\begin{table}[H]
\centering
\begin{tabular}{p{3.5cm}p{4cm}p{4cm}}
    \toprule
    \textbf{Facteur} & \textbf{Principe} & \textbf{Exemples} \\
    \midrule
    \textbf{Ce que tu connais} & Secret mémorisé. & Mot de passe, code PIN, phrase secrète. \\
    \textbf{Ce que tu possèdes} & Objet physique en possession. & Token matériel (RSA SecurID), carte à puce, OTP par SMS. \\
    \textbf{Ce que tu es} & Caractéristique biologique unique. & Empreinte digitale, iris, reconnaissance faciale. \\
    \textbf{Authentification Multi-Facteur (MFA)} & Combinaison de 2 facteurs ou plus. & Mot de passe + OTP SMS (ex. : Google Authenticator). \\
    \bottomrule
\end{tabular}
\caption{Facteurs d'authentification}
\end{table}

\begin{boxalert}
\textbf{Bonne pratique :} Le MFA est aujourd'hui indispensable pour tous les accès sensibles (VPN, messagerie professionnelle, accès administrateur). Un simple mot de passe n'est plus suffisant face aux attaques de phishing et de credential stuffing.
\end{boxalert}

% ------------------------------------------------------------------------------
\section{Implémentation des Contrôles d'Accès}
% ------------------------------------------------------------------------------

Les contrôles d'accès peuvent être implantés à trois niveaux :
\begin{itemize}
    \item \textbf{Physique :} Serrures, badges, caméras, sas de sécurité.
    \item \textbf{Logique :} ACL, pare-feu, chiffrement, comptes utilisateurs, MFA.
    \item \textbf{Administratif :} Politiques, procédures, formation, séparation des fonctions.
\end{itemize}

Le choix des contrôles doit être proportionné au niveau de risque et au budget de l'organisation (\textit{risk-based approach}). L'utilisation du \textbf{SSO (Single Sign-On)} simplifie la gestion des identités tout en maintenant un niveau de sécurité élevé.

% ------------------------------------------------------------------------------
\section{Modèles de Contrôles d'Accès (Rappel)}
% ------------------------------------------------------------------------------

Les modèles MAC, DAC, RBAC et ABAC ont été détaillés au Livre~III. Pour la gestion des identités, on retient particulièrement :
\begin{itemize}
    \item \textbf{RBAC} comme standard de facto en entreprise : un utilisateur hérite des droits de son rôle (Comptable, Admin, Auditeur, etc.).
    \item \textbf{Principe du moindre privilège} : n'accorder que les droits strictement nécessaires à la tâche.
    \item \textbf{Revue périodique des accès} : révocation rapide des droits lors de départ ou changement de poste.
\end{itemize}

% ------------------------------------------------------------------------------
\section{Évaluation des Acquis --- Livre~VI}
% ------------------------------------------------------------------------------

\begin{boxexo}
\textbf{QCM :}
\begin{enumerate}
    \item Qu'est-ce que l'IAAA ? \textbf{(a) Un protocole réseau \quad (b) Un modèle de contrôle d'accès \quad (c) Identification, Authentification, Autorisation, Accountability \quad (d) Un algorithme de chiffrement}
    \item L'authentification à deux facteurs combine : \textbf{(a) Deux mots de passe \quad (b) Un mot de passe et une empreinte digitale \quad (c) Deux cartes à puce \quad (d) Deux adresses email}
    \item Le principe du moindre privilège consiste à : \textbf{(a) Donner le maximum de droits \quad (b) N'accorder que les droits strictement nécessaires \quad (c) Partager les comptes entre collègues \quad (d) Désactiver l'authentification}
    \item SSO signifie : \textbf{(a) Secure Socket Option \quad (b) Single Sign-On \quad (c) System Security Object \quad (d) Secure Shell Operations}
\end{enumerate}

\textbf{Travail Pratique :}\\
Configurer un service d'authentification LDAP avec des groupes RBAC (Admin, Utilisateur, Auditeur) et tester les niveaux d'accès sur un répertoire partagé.
\end{boxexo}


% -------------------------------------------------------------------
% PARTIE VII — Évaluation et Tests de la Sécurité
% -------------------------------------------------------------------
\part{Évaluation et Tests de la Sécurité}

% Chapitre 7 — Tests d'intrusion et évaluation
% ==============================================================================
% Fichier : chapitres/07_evaluation_tests.tex
% Livre VII : Évaluation et Tests de la Sécurité
% ==============================================================================

\chapter{Évaluation et Tests de la Sécurité}
\label{chap:tests_securite}

\begin{boxobjectifs}
\textbf{Objectifs du Livre~VII :} Comprendre l'importance des tests de sécurité, maîtriser leurs types (boîte noire/grise/blanche) et les différentes étapes d'un test d'intrusion (autorisation, reconnaissance, scanning, exploitation, élévation, nettoyage, reporting).
\end{boxobjectifs}

% ------------------------------------------------------------------------------
\section{L'Importance des Tests de Sécurité}
% ------------------------------------------------------------------------------

Les tests de sécurité permettent d'identifier proactivement les vulnérabilités avant qu'un attaquant malveillant ne les exploite. Ils constituent un pilier fondamental de la démarche de sécurité défensive et sont souvent requis par les référentiels (ISO~27001, PCI~DSS).

% ------------------------------------------------------------------------------
\section{Types de Tests de Sécurité}
% ------------------------------------------------------------------------------

\begin{table}[H]
\centering
\small
\begin{tabular}{p{3cm}p{5.5cm}p{4cm}}
    \toprule
    \textbf{Type} & \textbf{Description} & \textbf{Usage} \\
    \midrule
    \textbf{Boîte Noire} & L'auditeur n'a aucune information sur la cible. Simule un attaquant externe. & Évaluation de l'exposition externe. \\
    \textbf{Boîte Grise}  & L'auditeur dispose d'informations partielles (ex. : compte standard). & Simulation d'un utilisateur interne. \\
    \textbf{Boîte Blanche} & Accès complet au code source et à l'architecture. & Audit de code, revue de sécurité approfondie. \\
    \textbf{Red Team}     & Équipe attaquante simulant une APT (menace persistante avancée). & Test de la capacité de détection et de réponse. \\
    \textbf{Blue Team}    & Équipe défensive (SOC) répondant aux attaques de la Red Team. & Amélioration des défenses et de la réactivité. \\
    \bottomrule
\end{tabular}
\caption{Types de tests de sécurité}
\end{table}

% ------------------------------------------------------------------------------
\section{Les Étapes d'un Test de Sécurité (Pentest)}
% ------------------------------------------------------------------------------

\subsection{a) Autorisation et Périmètre}
Étape préalable indispensable. Sans autorisation écrite (ROE — Rules of Engagement), toute intrusion est illégale. Définir : les systèmes inclus, les systèmes exclus, les actions autorisées, les horaires, les contacts en cas d'incident.

\subsection{b) Reconnaissance (Passive)}
Collecte d'informations sans contact direct avec la cible. Outils OSINT :
\begin{itemize}
    \item \textbf{Maltego :} Cartographie des relations (domaines, emails, organisations).
    \item \textbf{Shodan :} Moteur de recherche d'équipements connectés à Internet.
    \item \textbf{theHarvester :} Collecte d'emails, sous-domaines, IPs via moteurs de recherche.
    \item \textbf{WHOIS, dig, nslookup :} Informations DNS et d'enregistrement de domaine.
\end{itemize}

\subsection{c) Scanning et Énumération (Active)}
Interaction directe avec la cible pour cartographier les ports, services et systèmes.
\begin{itemize}
    \item \textbf{Nmap :} Scanner de ports et de systèmes d'exploitation.
    \item \textbf{Nessus / OpenVAS :} Scanners de vulnérabilités automatisés.
    \item \textbf{Angry IP Scanner, Advanced IP Scanner :} Découverte d'hôtes actifs sur un LAN.
    \item \textbf{Dirb / Gobuster :} Découverte de répertoires cachés sur les serveurs web.
\end{itemize}

\begin{lstlisting}[language=bash, caption=Commandes Nmap courantes]
# Scan de ports SYN (furtif)
nmap -sS -p 1-1024 192.168.1.0/24
# Détection de version et OS
nmap -sV -O 192.168.1.50
# Scan complet avec scripts par défaut
nmap -A 192.168.1.50
\end{lstlisting}

\subsection{d) Recherche des Vulnérabilités}
Analyse des résultats de scanning pour identifier les failles exploitables.
\begin{itemize}
    \item \textbf{CVE (Common Vulnerabilities and Exposures) :} Base de données des vulnérabilités publiques.
    \item \textbf{ExploitDB / SearchSploit :} Base d'exploits publics.
    \item \textbf{Rapid7 / Metasploit :} Référentiel d'exploits fonctionnels.
    \item \textbf{Nikto :} Scanner de vulnérabilités des serveurs web.
    \item \textbf{CVSS (Common Vulnerability Scoring System) :} Score de gravité de 0 à 10.
\end{itemize}

\subsection{e) L'Exploitation}
Tentative de compromission effective via les vulnérabilités identifiées.
\begin{itemize}
    \item \textbf{Metasploit Framework :} Standard de l'exploitation de vulnérabilités système.
    \item \textbf{BurpSuite / OWASP ZAP :} Proxy d'interception et d'attaque des applications web.
    \item \textbf{SQLMap :} Automatisation des injections SQL.
    \item \textbf{Mimikatz :} Extraction de credentials depuis la mémoire Windows (LSASS).
\end{itemize}

\subsection{f) L'Élévation de Privilèges}
Après l'accès initial, l'auditeur cherche à obtenir des droits administrateur/root pour prouver l'impact maximal.
\begin{itemize}
    \item \textbf{Linux :} Exploitation SUID, sudo misconfiguration, kernel exploits.
    \item \textbf{Windows :} Token impersonation, unquoted service path, AlwaysInstallElevated.
\end{itemize}

\subsection{g) Le Nettoyage}
Suppression de toutes les traces laissées (fichiers déposés, comptes créés, règles pare-feu modifiées) pour remettre le système dans son état initial.

\subsection{h) Le Reporting}
Rédaction du rapport de pentest incluant :
\begin{enumerate}
    \item \textbf{Synthèse exécutive} (pour la direction).
    \item \textbf{Méthodologie} et périmètre testé.
    \item \textbf{Findings} détaillés (titre, score CVSS, description, PoC, impact, remédiation).
    \item \textbf{Plan d'action} priorisé.
\end{enumerate}

% ------------------------------------------------------------------------------
\section{Évaluation des Acquis --- Livre~VII}
% ------------------------------------------------------------------------------

\begin{boxexo}
\textbf{QCM :}
\begin{enumerate}
    \item Un test boîte noire simule : \textbf{(a) Un auditeur interne avec accès complet \quad (b) Un attaquant externe sans information préalable \quad (c) Un utilisateur avec un compte standard \quad (d) Un développeur avec le code source}
    \item Nmap est un outil de : \textbf{(a) Chiffrement \quad (b) Scanning de ports \quad (c) Gestion des identités \quad (d) Journalisation}
    \item Le CVE est : \textbf{(a) Un outil d'exploitation \quad (b) Un dictionnaire public des vulnérabilités \quad (c) Un protocole réseau \quad (d) Un modèle de contrôle d'accès}
    \item Le nettoyage en pentest consiste à : \textbf{(a) Formater les disques \quad (b) Supprimer les traces laissées pendant le test \quad (c) Réinstaller le système \quad (d) Désactiver les pare-feux}
\end{enumerate}

\textbf{Travail Pratique (ENSPY) :}\\
Réaliser une évaluation de la sécurité de l'application de gestion des notes de l'ENSPY. Identifier les vulnérabilités (OWASP Top 10), les scorer avec CVSS et proposer des remédiations.
\end{boxexo}


% -------------------------------------------------------------------
% PARTIE VIII — Opérations de Sécurité
% -------------------------------------------------------------------
\part{Opérations de Sécurité}

% Chapitre 8 — Gestion des opérations de sécurité
% ==============================================================================
% Fichier : chapitres/08_operations_securite.tex
% Livre VIII : Opérations de Sécurité
% ==============================================================================

\chapter{Opérations de Sécurité}
\label{chap:operations}

\begin{boxobjectifs}
\textbf{Objectifs du Livre~VIII :} Présenter les opérations quotidiennes de sécurité informatique : gouvernance administrative, investigation numérique, gestion des incidents, contrôles opérationnels, tolérance aux pannes et plans de continuité/reprise d'activité.
\end{boxobjectifs}

% ------------------------------------------------------------------------------
\section{Sécurité Administrative}
% ------------------------------------------------------------------------------

\subsection{Gouvernance de la Sécurité}
La gouvernance de la sécurité est le processus de gestion et de supervision de la sécurité informatique d'une entreprise. Elle définit les objectifs, les rôles, les responsabilités et les processus de décision en matière de sécurité, en les alignant sur la stratégie générale de l'organisation.

\subsection{Politiques de Sécurité}
Les politiques de sécurité décrivent les règles et procédures à respecter. Elles doivent être :
\begin{itemize}
    \item Validées par les autorités dirigeantes.
    \item Claires, concises et accessibles à tous les employés.
    \item Régulièrement mises à jour face à l'évolution des menaces.
\end{itemize}

\subsection{Sensibilisation à la Sécurité Informatique}
La composante humaine est le maillon le plus faible de la sécurité. La sensibilisation peut prendre la forme de :
\begin{itemize}
    \item Formations périodiques (e-learning, ateliers).
    \item Campagnes de phishing simulé pour tester la vigilance.
    \item Newsletters et affiches de rappel des bonnes pratiques.
\end{itemize}

\subsection{Conformité Réglementaire}
L'organisation doit se conformer aux lois et réglementations applicables (RGPD, loi camerounaise sur la cybersécurité, PCI~DSS, ISO~27001) sous peine de sanctions administratives et pénales.

\subsection{Audits de Sécurité}
Les audits internes et externes évaluent l'efficacité des contrôles en place et leur conformité aux politiques. Ils peuvent être : audits de vulnérabilité, de conformité, de configuration, de code ou tests d'intrusion.

% ------------------------------------------------------------------------------
\section{Investigation Numérique (Digital Forensics)}
% ------------------------------------------------------------------------------

\begin{boxdef}
\textbf{Investigation numérique :} Discipline consistant à collecter, préserver, analyser et présenter des preuves numériques de manière légalement admissible, dans le but d'établir les faits liés à un incident de sécurité ou à une infraction.
\end{boxdef}

\subsection{Chaîne de Garde (Chain of Custody)}
La chaîne de garde assure la traçabilité et l'intégrité des preuves numériques depuis leur collecte jusqu'à leur présentation devant un tribunal. Elle implique :
\begin{itemize}
    \item Enregistrement de toutes les actions effectuées sur les preuves (qui, quoi, quand, pourquoi).
    \item Scellage et stockage sécurisé des originaux.
    \item Travail exclusivement sur des copies forensiques.
\end{itemize}

\subsection{Collecte des Preuves}
\begin{itemize}
    \item \textbf{Ne jamais travailler sur l'original} : créer une copie bit-à-bit (\textit{forensic image}).
    \item \textbf{Données volatiles en priorité} : RAM, processus actifs, connexions réseau ouvertes (disparaissent à l'extinction).
    \item \textbf{Outils spécialisés :} FTK~Imager, EnCase, dd (Linux), Volatility (mémoire vive).
\end{itemize}

\begin{lstlisting}[language=bash, caption=Création d'une image forensique avec dd]
# Créer une copie bit-à-bit d'un disque
dd if=/dev/sdb of=/mnt/evidence/image.dd bs=4096 conv=noerror,sync
# Calculer et vérifier le hash MD5 pour l'intégrité
md5sum /dev/sdb > hash_original.md5
md5sum /mnt/evidence/image.dd > hash_copie.md5
diff hash_original.md5 hash_copie.md5   # Doit être identique
\end{lstlisting}

\subsection{Analyse des Éléments Collectés}
L'analyse porte sur :
\begin{itemize}
    \item Journaux système (Windows Event Logs, syslog Linux, journaux pare-feu).
    \item Artefacts navigateur (historique, cache, cookies).
    \item Fichiers supprimés récupérés depuis l'espace non alloué.
    \item Trafic réseau capturé (pcap analysé avec Wireshark).
    \item Artefacts mémoire (processus cachés, connexions actives, mots de passe en mémoire).
\end{itemize}

\subsection{Rédaction du Rapport d'Investigation}
Le rapport doit être factuel, précis et compréhensible par des non-techniciens (juges, avocats). Il contient :
\begin{enumerate}
    \item Contexte et mandat de la mission.
    \item Méthodologie et outils utilisés.
    \item Chronologie des événements reconstitués.
    \item Preuves collectées et leur signification.
    \item Conclusions et recommandations.
\end{enumerate}

% ------------------------------------------------------------------------------
\section{Gestion des Incidents}
% ------------------------------------------------------------------------------

\begin{center}
\begin{tikzpicture}[
    node distance=1.4cm,
    box/.style={rectangle, rounded corners=4pt, draw=CouleurPrimaire,
                fill=CouleurPrimaire!10, text width=3.2cm, align=center,
                minimum height=1.1cm, font=\small\bfseries}
]
    \node[box] (prep)   {1. Préparation};
    \node[box] (det)    [right=of prep]  {2. Détection \& Analyse};
    \node[box] (cont)   [right=of det]   {3. Confinement \& Éradication};
    \node[box] (recov)  [below=of cont]  {4. Recouvrement};
    \node[box] (post)   [left=of recov]  {5. Activités Post-incident};
    \node[box] (learn)  [left=of post]   {6. Leçons tirées};

    \draw[->, thick, CouleurPrimaire] (prep) -- (det);
    \draw[->, thick, CouleurPrimaire] (det) -- (cont);
    \draw[->, thick, CouleurPrimaire] (cont) -- (recov);
    \draw[->, thick, CouleurPrimaire] (recov) -- (post);
    \draw[->, thick, CouleurPrimaire] (post) -- (learn);
    \draw[->, thick, CouleurPrimaire, dashed] (learn) -- (prep);
\end{tikzpicture}
\captionof{figure}{Cycle de gestion des incidents (NIST SP 800-61)}
\end{center}

\subsection{Préparation}
Mettre en place les outils, les équipes, les procédures et les communications nécessaires \textit{avant} qu'un incident ne survienne :
\begin{itemize}
    \item Constitution d'une équipe CSIRT (Computer Security Incident Response Team).
    \item Définition des niveaux de gravité et des escalades.
    \item Déploiement d'un SIEM (Security Information and Event Management).
\end{itemize}

\subsection{Détection et Analyse}
Identifier et qualifier l'incident :
\begin{itemize}
    \item Sources de détection : alertes IDS/IPS, SIEM, antivirus, signalement utilisateur.
    \item Triage : est-ce un faux positif ou un incident réel ?
    \item Qualification de la gravité (P1~Critique à P4~Faible).
\end{itemize}

\subsection{Confinement, Éradication et Recouvrement}
\begin{itemize}
    \item \textbf{Confinement court terme :} Isolation du système compromis du réseau.
    \item \textbf{Confinement long terme :} Patch, changement de credentials.
    \item \textbf{Éradication :} Suppression du malware, fermeture des accès non autorisés.
    \item \textbf{Recouvrement :} Restauration depuis une sauvegarde saine, remise en production surveillée.
\end{itemize}

\subsection{Activités Post-incident}
\begin{itemize}
    \item Rédaction d'un rapport post-mortem.
    \item Identification des causes racines (Root Cause Analysis).
    \item Mise à jour des politiques et procédures.
    \item Communication aux autorités si requis (RGPD : notification sous 72~h).
\end{itemize}

% ------------------------------------------------------------------------------
\section{Contrôles Opérationnels}
% ------------------------------------------------------------------------------

\begin{table}[H]
\centering
\small
\begin{tabular}{p{3cm}p{9cm}}
    \toprule
    \textbf{Type} & \textbf{Exemples} \\
    \midrule
    \textbf{Contrôles physiques}       & Serrures, cages serveurs, caméras, badges d'accès, agents de sécurité. \\
    \textbf{Contrôles logiques}        & Pare-feux, IDS/IPS, chiffrement, antivirus, journalisation, MFA. \\
    \textbf{Contrôles administratifs}  & Politiques, procédures, formations, séparation des fonctions, audits. \\
    \bottomrule
\end{tabular}
\caption{Catégories de contrôles opérationnels}
\end{table}

% ------------------------------------------------------------------------------
\section{Tolérance aux Pannes}
% ------------------------------------------------------------------------------

\subsection{Haute Disponibilité (HA)}
Conception d'un système pour minimiser les interruptions non planifiées. Techniques :
\begin{itemize}
    \item \textbf{Clustering :} Plusieurs serveurs actifs simultanément (actif-actif).
    \item \textbf{Failover :} Bascule automatique vers un serveur de secours (actif-passif).
    \item \textbf{Load Balancing :} Répartition de la charge entre plusieurs serveurs.
\end{itemize}

\subsection{Architectures Redondantes}
\begin{itemize}
    \item \textbf{RAID (Redundant Array of Independent Disks) :} Redondance des données sur plusieurs disques.
    \item \textbf{Double alimentation électrique + UPS :} Protection contre les coupures de courant.
    \item \textbf{Double connectivité réseau (liaisons redondantes) :} Continuité en cas de panne d'un opérateur.
\end{itemize}

% ------------------------------------------------------------------------------
\section{Plan de Continuité des Activités (PCA) et Plan de Reprise (PRA)}
% ------------------------------------------------------------------------------

\begin{boxdef}
\textbf{PCA (Plan de Continuité des Activités) :} Ensemble de mesures permettant de maintenir les fonctions critiques à un niveau acceptable \textit{pendant} un sinistre.

\textbf{PRA (Plan de Reprise après Sinistre) :} Ensemble de mesures permettant de reprendre l'activité normale \textit{après} un sinistre. Il définit les procédures de retour à l'état nominal.
\end{boxdef}

\begin{table}[H]
\centering
\small
\begin{tabular}{p{3cm}p{8.5cm}}
    \toprule
    \textbf{Indicateur} & \textbf{Définition} \\
    \midrule
    \textbf{RTO} (Recovery Time Objective)  & Durée maximale acceptable d'interruption. Ex. : 4~heures. \\
    \textbf{RPO} (Recovery Point Objective) & Perte de données maximale acceptable. Ex. : données des 2 dernières heures. \\
    \textbf{MTD} (Maximum Tolerable Downtime) & Durée maximale au-delà de laquelle l'organisation ne peut plus fonctionner. \\
    \bottomrule
\end{tabular}
\caption{Indicateurs clés de continuité}
\end{table}

% ------------------------------------------------------------------------------
\section{Évaluation des Acquis --- Livre~VIII}
% ------------------------------------------------------------------------------

\begin{boxexo}
\textbf{QCM :}
\begin{enumerate}
    \item La chaîne de garde (chain of custody) vise principalement à : \textbf{(a) Accélérer l'enquête \quad (b) Assurer la traçabilité et l'intégrité des preuves \quad (c) Supprimer les données compromises \quad (d) Former les enquêteurs}
    \item Dans le cycle NIST de gestion des incidents, quelle phase suit immédiatement la détection ? \textbf{(a) Préparation \quad (b) Confinement \quad (c) Recouvrement \quad (d) Rédaction du rapport}
    \item Le RPO (Recovery Point Objective) mesure : \textbf{(a) Le temps de reprise maximal \quad (b) La perte de données maximale acceptable \quad (c) Le coût de la reprise \quad (d) La disponibilité du système}
    \item Le SIEM est un outil de : \textbf{(a) Chiffrement des données \quad (b) Collecte et corrélation des journaux de sécurité \quad (c) Gestion des identités \quad (d) Sauvegarde automatique}
\end{enumerate}

\textbf{Travaux Pratiques :}
\begin{itemize}
    \item Cloner un disque avec FTK~Imager et vérifier l'intégrité par hash MD5/SHA-256.
    \item Prendre en main l'outil EnCase pour l'analyse forensique d'une image disque.
    \item Proposer une campagne de sensibilisation aux cyberattaques pour l'ENSPY (affiche, quiz, simulation de phishing).
\end{itemize}
\end{boxexo}


% -------------------------------------------------------------------
% PARTIE IX — Sécurité des Applications
% -------------------------------------------------------------------
\part{Sécurité des Applications}

% Chapitre 9 — Développement sécurisé et DevSecOps
% ==============================================================================
% Fichier : chapitres/09_securite_applications.tex
% Livre IX : Sécurité des Applications
% ==============================================================================

\chapter{Sécurité des Applications}
\label{chap:securite_apps}

\begin{boxobjectifs}
\textbf{Objectifs du Livre~IX :} Maîtriser les méthodes de développement sécurisé (SSDLC), l'évaluation de la sécurité des logiciels (boîte noire/grise/blanche), les bases de données (SQL, NoSQL, injections), la programmation orientée objet sécurisée et les principes du code défensif.
\end{boxobjectifs}

% ------------------------------------------------------------------------------
\section{Méthodes de Développement}
% ------------------------------------------------------------------------------

\subsection{Méthodes Classiques}
Les méthodes séquentielles (cascade, modèle en~V, itératif) imposent une planification détaillée et une documentation complète à chaque phase. Elles manquent de flexibilité face aux changements tardifs d'exigences, notamment les exigences de sécurité découvertes en fin de cycle.

\subsection{Méthodes Agiles}
Scrum, XP (Extreme Programming) et Kanban privilégient des itérations courtes (sprints), des livraisons fréquentes et l'adaptation continue. Elles facilitent l'intégration de la sécurité à chaque sprint (\textit{Security by Design}).

\subsection{SSDLC (Secure Software Development Life Cycle)}
\begin{boxdef}
Le \textbf{SSDLC} est un cycle de développement logiciel qui intègre la sécurité à \textit{chaque étape} : de la définition des exigences à la mise en production et à la maintenance. L'objectif est de «~déplacer la sécurité vers la gauche~» (\textit{Shift Left Security}) pour détecter les vulnérabilités le plus tôt possible, là où leur correction coûte le moins cher.
\end{boxdef}

\begin{table}[H]
\centering
\small
\begin{tabular}{p{3.5cm}p{8cm}}
    \toprule
    \textbf{Phase SSDLC} & \textbf{Activité de Sécurité} \\
    \midrule
    Exigences         & Définir les exigences de sécurité (authentication, autorisation, chiffrement). \\
    Conception        & Modélisation des menaces (STRIDE, DREAD), architecture sécurisée. \\
    Développement     & Revue de code sécurisé, respect des bonnes pratiques (OWASP). \\
    Test              & Tests de sécurité (SAST, DAST, pentest, fuzzing). \\
    Déploiement       & Configuration sécurisée, scan de vulnérabilités pré-production. \\
    Maintenance       & Gestion des patchs, surveillance continue, réponse aux incidents. \\
    \bottomrule
\end{tabular}
\caption{Activités de sécurité par phase SSDLC}
\end{table}

% ------------------------------------------------------------------------------
\section{Évaluation de la Sécurité des Logiciels}
% ------------------------------------------------------------------------------

\subsection{Vulnérabilités Logicielles Communes}

\begin{table}[H]
\centering
\small
\begin{tabular}{p{4cm}p{8cm}}
    \toprule
    \textbf{Vulnérabilité} & \textbf{Description} \\
    \midrule
    \textbf{Injection de code} & Exploitation d'une entrée utilisateur pour injecter du code (SQL, shell, JavaScript) dans l'application. \\
    \textbf{Exposition de données sensibles} & Exposition involontaire de mots de passe, clés d'API, informations personnelles. \\
    \textbf{Dépassement de tampon (Buffer Overflow)} & Écriture au-delà de la mémoire allouée permettant l'exécution de code arbitraire. \\
    \textbf{Vulnérabilités des composants tiers} & Bibliothèques/frameworks présentant des CVE connus non patchés. \\
    \textbf{Manque de validation des entrées} & Absence de filtrage des données saisies par l'utilisateur. \\
    \textbf{Mauvaise gestion des erreurs} & Messages d'erreur trop verbeux révélant la technologie ou l'architecture interne. \\
    \bottomrule
\end{tabular}
\caption{Vulnérabilités logicielles courantes}
\end{table}

\subsection{Types de Tests de Sécurité Logicielle}

\begin{itemize}
    \item \textbf{SAST (Static Application Security Testing) :} Analyse du code source sans exécution. Détecte les vulnérabilités tôt dans le cycle.
    \item \textbf{DAST (Dynamic Application Security Testing) :} Tests en boîte noire sur l'application en cours d'exécution (ex. : BurpSuite, OWASP ZAP).
    \item \textbf{IAST (Interactive AST) :} Combinaison de SAST et DAST via des agents instrumentant l'application pendant les tests.
    \item \textbf{Fuzzing :} Envoi de données aléatoires ou malformées pour provoquer des comportements inattendus.
\end{itemize}

\subsection{Méthodologie de Test de Sécurité Applicative}
\begin{enumerate}
    \item \textbf{Reconnaissance :} Identification de la technologie, des endpoints et de l'architecture (Maltego, Shodan, Nmap, Burp Spider).
    \item \textbf{Recherche des vulnérabilités :} OWASP~ZAP, SQLMap, DirBuster, Nikto.
    \item \textbf{Exploitation :} Metasploit, BurpSuite, BeEF (XSS).
    \item \textbf{Élévation de privilèges :} Mimikatz, PowerSploit, Empire.
    \item \textbf{Reporting :} Fiches CVSS, impact business, recommandations de correction.
\end{enumerate}

% ------------------------------------------------------------------------------
\section{Sécurité des Bases de Données}
% ------------------------------------------------------------------------------

\subsection{Bases de Données SQL}
Les SGBD relationnels (MySQL, PostgreSQL, Oracle, SQL Server, MariaDB) organisent les données en tables liées par des clés primaires et étrangères. Les données sont interrogées via le langage SQL.

\begin{boxalert}
\textbf{Injection SQL :} Vulnérabilité la plus critique des applications web (OWASP \#1). Elle survient lorsque des entrées utilisateur non filtrées sont incorporées dans une requête SQL.

\begin{lstlisting}[language=SQL]
-- Requête vulnérable
SELECT * FROM users WHERE username='$input' AND password='$pass';

-- Payload d'attaque : ' OR '1'='1' --
-- La requête devient :
SELECT * FROM users WHERE username='' OR '1'='1' --' AND password='...';
-- Résultat : toujours vraie -> accès accordé sans mot de passe
\end{lstlisting}

\textbf{Contre-mesure :} Toujours utiliser des \textbf{requêtes préparées (prepared statements)} :
\begin{lstlisting}[language=PHP]
$stmt = $pdo->prepare(
    "SELECT * FROM users WHERE username = ? AND password = ?"
);
$stmt->execute([$username, $hashed_password]);
\end{lstlisting}
\end{boxalert}

\subsection{Bases de Données NoSQL}
Les SGBD non relationnels (MongoDB, Redis, Cassandra, CouchBase) utilisent des modèles de données flexibles :

\begin{table}[H]
\centering
\small
\begin{tabular}{p{3cm}p{4cm}p{4.5cm}}
    \toprule
    \textbf{Type NoSQL} & \textbf{Modèle de données} & \textbf{Exemples de SGBD} \\
    \midrule
    Clé-valeur       & Paires clé $\Rightarrow$ valeur simples. & Redis, Riak \\
    Orienté document & Documents JSON / XML. & MongoDB, CouchBase \\
    Orienté colonnes & Familles de colonnes. & Cassandra, HBase \\
    Graphe           & Nœuds et relations. & Neo4j, Amazon Neptune \\
    \bottomrule
\end{tabular}
\caption{Types de bases de données NoSQL}
\end{table}

\subsection{Injections NoSQL}
Les SGBD NoSQL ne sont pas immunisés contre les injections. Exemple MongoDB :
\begin{lstlisting}[language=javascript, caption=Injection NoSQL MongoDB]
// Requête vulnérable (Node.js / Express)
db.users.find({ username: req.body.user, password: req.body.pass });

// Payload : { "$gt": "" }  -> condition toujours vraie
// username: { "$gt": "" }, password: { "$gt": "" }  => accès accordé
\end{lstlisting}

% ------------------------------------------------------------------------------
\section{Programmation Orientée Objet (POO) et Sécurité}
% ------------------------------------------------------------------------------

\begin{table}[H]
\centering
\small
\begin{tabular}{p{3cm}p{8.5cm}}
    \toprule
    \textbf{Principe POO} & \textbf{Intérêt pour la sécurité} \\
    \midrule
    \textbf{Encapsulation}  & Masque les détails d'implémentation. Protège les attributs sensibles (private/protected). \\
    \textbf{Héritage}       & Réutilisation des contrôles de sécurité dans les classes dérivées. Risque : héritage de vulnérabilités. \\
    \textbf{Polymorphisme}  & Permet de substituer des implémentations de chiffrement sans changer l'interface. \\
    \textbf{Abstraction}    & Interfaces et classes abstraites masquant la complexité cryptographique. \\
    \textbf{Couplage faible} & Composants indépendants : une vulnérabilité dans un module n'expose pas les autres. \\
    \textbf{Cohésion forte} & Un module fait une seule chose bien : facilite la revue et le test de sécurité. \\
    \bottomrule
\end{tabular}
\caption{Principes POO appliqués à la sécurité}
\end{table}

% ------------------------------------------------------------------------------
\section{Programmation Sécurisée}
% ------------------------------------------------------------------------------

\begin{boxinfo}
\textbf{Principes du code défensif (OWASP Secure Coding Practices) :}
\begin{itemize}
    \item \textbf{Valider toutes les entrées :} Whitelist (liste des valeurs acceptées), jamais de confiance implicite.
    \item \textbf{Encoder les sorties :} Prévenir les attaques XSS par échappement HTML/JS.
    \item \textbf{Authentification forte :} Stocker les mots de passe avec bcrypt/Argon2 (jamais MD5 ou SHA-1 seul).
    \item \textbf{Gestion des sessions :} Tokens aléatoires et suffisamment longs, expiration, régénération après login.
    \item \textbf{Contrôle d'accès :} Vérifier les autorisations côté serveur, pas seulement côté client.
    \item \textbf{Cryptographie :} Utiliser les bibliothèques standard (OpenSSL, libsodium), ne jamais implémenter ses propres algorithmes.
    \item \textbf{Gestion des erreurs :} Ne jamais exposer de détails techniques dans les messages d'erreur utilisateur.
    \item \textbf{Journalisation :} Enregistrer les événements de sécurité (tentatives de connexion, accès refusés) sans consigner de données sensibles.
    \item \textbf{Sécurité des composants tiers :} Tenir à jour les dépendances (npm audit, pip check, OWASP Dependency-Check).
\end{itemize}
\end{boxinfo}

% ------------------------------------------------------------------------------
\section{Évaluation des Acquis --- Livre~IX}
% ------------------------------------------------------------------------------

\begin{boxexo}
\textbf{QCM :}
\begin{enumerate}
    \item Le SSDLC est un cycle de développement qui : \textbf{(a) Accélère le développement \quad (b) Intègre la sécurité à chaque étape \quad (c) Se concentre sur la performance \quad (d) Élimine les tests}
    \item Une vulnérabilité logicielle est : \textbf{(a) Une fonctionnalité utile \quad (b) Une faiblesse pouvant être exploitée \quad (c) Une erreur corrigée par un patch \quad (d) Un composant tiers}
    \item La meilleure contre-mesure contre les injections SQL est : \textbf{(a) Utiliser NoSQL \quad (b) Désactiver les logs \quad (c) Utiliser des requêtes préparées \quad (d) Chiffrer la base de données}
    \item L'encapsulation en POO contribue à la sécurité en : \textbf{(a) Accélérant l'exécution \quad (b) Masquant les attributs et implémentations sensibles \quad (c) Permettant l'héritage multiple \quad (d) Simplifiant l'interface}
    \item Le DAST (Dynamic Application Security Testing) teste : \textbf{(a) Le code source statique \quad (b) L'application en cours d'exécution \quad (c) La base de données uniquement \quad (d) L'infrastructure réseau}
    \item Le principe «~Shift Left Security~» signifie : \textbf{(a) Déplacer les tests en fin de projet \quad (b) Intégrer la sécurité dès les premières phases du développement \quad (c) Sécuriser uniquement l'interface utilisateur \quad (d) Réduire le budget sécurité}
\end{enumerate}

\textbf{Travail Pratique :}\\
Réaliser une évaluation complète de la sécurité de l'application de gestion des notes de l'ENSPY :
\begin{enumerate}
    \item Identifier les vulnérabilités OWASP~Top~10 présentes.
    \item Tester l'injection SQL sur les formulaires de connexion et de recherche.
    \item Proposer et implémenter les corrections (requêtes préparées, validation des entrées, hachage bcrypt).
    \item Rédiger un rapport avec fiches CVSS pour chaque vulnérabilité identifiée.
\end{enumerate}
\end{boxexo}


% --- Conclusion Générale ---
\input{chapitres/00_conclusion}

% ==============================================================================
% MATIÈRE POSTLIMINAIRE
% ==============================================================================
\backmatter

% --- Bibliographie ---
\clearpage
\printbibliography[
    heading=bibintoc,
    title={Bibliographie et Références}
]

% --- À Propos de l'Auteur ---
\clearpage
\input{chapitres/Apropos_Auteur}

% --- Index ---
\clearpage
\printindex
\addcontentsline{toc}{chapter}{Index}

% ==============================================================================
% 4ÈME DE COUVERTURE (IDENTIQUE AU MODÈLE)
% ==============================================================================
\clearpage
\thispagestyle{empty}
\pagecolor{CouleurPrimaire}

\begin{tikzpicture}[remember picture, overlay]
    \fill[CouleurAccent]
        (current page.south west) rectangle
        ([yshift=1.8cm] current page.south east);
    \fill[CouleurAccent]
        ([yshift=-1.4cm] current page.north west) rectangle
        ([yshift=-1.7cm] current page.north east);
\end{tikzpicture}

\vspace*{2.5cm}

\begin{center}
    {\color{CouleurAccent}\rule{0.6\linewidth}{1pt}}\\[1.5ex]
    {\color{white}\Large\bfseries À propos de ce cours}\\[1.5ex]
    {\color{CouleurAccent}\rule{0.6\linewidth}{1pt}}
\end{center}

\vspace{1cm}

\begin{center}
\begin{minipage}{0.75\textwidth}
\color{white!85!gray}\setstretch{1.4}
Ce cours propose une approche intégrée et rigoureuse de la sécurité informatique, articulant les fondamentaux théoriques, les pratiques opérationnelles et les référentiels internationaux reconnus (ISO 27001, NIST, CISSP).

\medskip

Il s'adresse aux étudiants en informatique souhaitant acquérir les compétences conceptuelles et opérationnelles requises pour concevoir, mettre en œuvre et évaluer la sécurité des systèmes d'information avec professionnalisme.

\medskip

Les thèmes abordés couvrent la gestion des risques, la cryptographie, la sécurité des réseaux, la gestion des identités et des accès, les tests d'intrusion, la sécurité des applications et la gouvernance de la cybersécurité.
\end{minipage}
\end{center}

\vfill

\begin{center}
    {\color{CouleurAccent}\rule{0.7\linewidth}{1.2pt}}\\[1.2ex]
    {\color{white}\large\bfseries MINKA MI NGUIDJOÏ Thierry Emmanuel}\\[0.4ex]
    {\color{white!80!gray}\small
        CISA~·~CISM~·~CGEIT~·~CRISC~·~CDPSE~·~ISO 27001~·~ISO 22301~·~ISO 9001}\\[0.3ex]
    {\color{white!70!gray}\footnotesize\itshape
        Chercheur — LIMSI — ENSP, Université de Yaoundé~I}\\[0.3ex]
    {\color{white!70!gray}\footnotesize\itshape
        Sr-Eng~·~Sr-Auditor~·~GRC Expert~·~Cryptographe}\\[1.2ex]
    {\color{CouleurAccent}\rule{0.7\linewidth}{1.2pt}}
\end{center}

\vspace{1cm}

\afterpage{\nopagecolor}

\end{document}
% ==============================================================================
% Fin de Cours_Securite.tex
% ==============================================================================
